% !TeX program = xelatex
% Methodological framework with a targeted evidence review, 2026-10-05; final version 1.0.0.
% Author supplied: 游尹哲郡 / Youyin Zhejun; Central University of Finance and Economics.
% Single-file manuscript with embedded references.
% References are embedded. Numerical examples: scripts/reproduce_examples.py.
\documentclass[11pt,a4paper]{article}
\usepackage[T1]{fontenc}
\usepackage{lmodern}
\usepackage[margin=28mm,headheight=15pt]{geometry}
\usepackage{amsmath,amssymb,amsthm,booktabs,array,tabularx,longtable}
\usepackage[authoryear,round]{natbib}
\usepackage{xcolor,fancyhdr,enumitem,float}
\usepackage[section]{placeins}
\usepackage{xurl}
\usepackage[unicode,colorlinks=true,linkcolor=black,citecolor=black,urlcolor=blue!45!black]{hyperref}
\hypersetup{pdftitle={Semantic Decision Interfaces for Statistical Measurement and Population Estimation: A Methodological Framework},pdfauthor={Youyin Zhejun},pdfsubject={Methodological framework with a targeted evidence review},pdfkeywords={semantic decision interfaces,statistical measurement,population estimation,verification design,model-assisted inference}}
\urlstyle{same}
\linespread{1.08}
\setlength{\parindent}{1.5em}
\setlength{\parskip}{0.15em}
\setlength{\emergencystretch}{2em}
\setlength{\tabcolsep}{4pt}
\renewcommand{\arraystretch}{1.18}
\pagestyle{fancy}
\fancyhf{}
\fancyhead[L]{\small Youyin Zhejun}
\fancyhead[R]{\small Semantic decision interfaces: a methodological framework}
\fancyfoot[C]{\small\thepage}
\renewcommand{\headrulewidth}{0.3pt}
\setlist[enumerate]{leftmargin=2em,itemsep=0.25em,topsep=0.4em}
\newcolumntype{Y}{>{\raggedright\arraybackslash}X}
\newcolumntype{L}[1]{>{\raggedright\arraybackslash}p{#1}}
\newcommand{\E}{\mathbb{E}}
\newcommand{\Prob}{\mathbb{P}}
\newcommand{\Var}{\operatorname{Var}}
\newcommand{\Cov}{\operatorname{Cov}}
\newcommand{\ind}{\mathbf{1}}
\DeclareMathOperator*{\argmin}{arg\,min}
\DeclareMathOperator*{\argmax}{arg\,max}
\title{\bfseries Semantic Decision Interfaces for Statistical Measurement and Population Estimation:\\A Methodological Framework}
\author{Youyin Zhejun\\[0.35em]
\small Central University of Finance and Economics\\
\small ORCID: \href{https://orcid.org/0009-0009-5316-8193}{0009-0009-5316-8193}}
\date{5 October 2026}

\begin{document}
\maketitle
\thispagestyle{plain}
\begin{abstract}
Semantic decision interfaces evaluate records against explicit semantic criteria within prescribed answer spaces. This article develops a methodological framework for their use in statistical measurement and population estimation, supported by a targeted evidence review of nine studies concentrated on Jev and compatible implementations, together with interface documentation. The framework connects the estimand, measurement protocol, interface output, validation conditions, verification design, and estimator. It distinguishes the substantive object of study, the protocol-defined variable, and observed verification labels. Constructed examples show how these distinctions identify prevalence bias, incompatible pair judgments, coding-dependent ordinal summaries, and differences in inferential efficiency. Established difference estimators and allocation methods supply verification schemes for independent samples and fixed corpora. A sampling-design evaluation on an SMS corpus with a fixed natural language inference model compares five prevalence estimators. Pilot-based coefficient tuning improves on unit-weight two-stage correction, while direct label sampling has the lowest MSE among the verified options at both evaluated budgets. Three analytical principles organize the framework: semantic judgments participate in measurement; local prediction quality and downstream validity have distinct requirements; and statistical usefulness is evaluated for a specified estimand and verification budget. The evidence review illustrates implementation choices and identifies statistical uses that remain proposals for evaluation.
\end{abstract}
\noindent\textbf{Keywords:} semantic decision interfaces; statistical measurement; population estimation; verification design; probability calibration; model-assisted inference.

\section{Introduction}\label{sec:intro}

A study of customer complaints may estimate the prevalence of billing disputes, group reports by underlying service failure, or compare the severity of incidents across regions. Each analysis begins with a measurement decision: which features of a narrative count as evidence, which distinctions define a category, and which records describe the same event. The resulting variables subsequently enter familiar statistical operations, including tabulation, regression, and uncertainty estimation. Their construction is therefore part of the statistical model, even when the construction procedure is implemented through a software service.

Jev, introduced by TypeSafe AI as a System One model, accepts a state and typed questions and returns structured answers with probabilities \citep{typesafe-intro}. This article uses \emph{semantic decision interface} for an operation that evaluates supplied records against explicit semantic criteria within a prescribed answer space. A model implements that operation; its architecture, probability readout, and use of persistent context are implementation properties. System One denotes terminology used in the current model ecosystem. The statistical question is how the specified judgments support measurement and estimation.

Several established research traditions address the expensive acquisition of statistical variables. Graph-based semi-supervised learning propagates information from labeled observations through a similarity structure \citep{zhu2003}. Weak supervision combines noisy rules or annotation sources \citep{ratner2017}. Zero-shot classification uses label descriptions to specify a new classification task \citep{yin2019}. Human judgments can also constrain a clustering solution \citep{wagstaff2001}. These traditions supply methods for using the information elicited through a semantic decision interface. Their assumptions help identify which requested judgments support the intended statistical task.

The distinction between predicting individual labels and estimating a population quantity is especially consequential. \citet{hopkins2010} developed automated content analysis around category proportions, showing why an evaluation focused on individual classifications can miss the research objective. \citet{grimmer2013} emphasized substantive interpretation and application-specific validation in text analysis. Contemporary language-model evaluations continue this discussion: \citet{ziems2024} assessed zero-shot annotation across social-science tasks, and \citet{egami2023} studied the effect of imperfect surrogate labels on downstream inference. Decision models belong within this methodological history.

The neighboring reviews already cover substantial parts of this problem. \citet{tan2024} organize annotation generation, assessment, and utilization, including prompt construction, annotation criteria, and comparison with reference labels. \citet{zhanggao2026} organize prediction-powered inference around rectification, label acquisition, predictor dependence, and validity targets, and relate it to measurement-error methods. Their baseline treats $Y$ as the gold-standard outcome. The present framework makes four connected choices explicit: the coding protocol and its observed verification variable; the information supplied to a query and its answer space; structural consistency across related queries; and the resulting constraints on verification and estimation. It separates protocol implementation from substantive measurement, treats input and answer-space changes as changes to the measurement operation, and links pairwise or hierarchical judgments to joint consistency requirements. These choices determine which reference labels, sampling frame, residuals, and variance calculations an analysis requires. The increment is this organized connection between measurement choices and established inferential tools. Section~\ref{sec:inference} applies those tools; a real-corpus evaluation examines their behavior under a fixed prediction rule and labeling budget.

A targeted evidence review of nine directly relevant studies describes an early empirical literature concentrated on Jev and compatible implementations. General statistical theory supplies the estimation arguments, interface specifications define the operations, and implementation studies supply evidence on selected tasks. Constructed examples and the corpus evaluation demonstrate the consequences of particular measurement and sampling choices.

The article makes three contributions. First, it specifies how semantic judgments enter measurement, distinguishing the substantive object, the protocol-defined variable, and the labels obtained through verification. Second, it maps estimands to variable construction and validation conditions, so that researchers can select checks for prevalence, partitions, ordinal summaries, and linked units. Third, it organizes verification and estimation schemes by sampling conditions, relating auxiliary predictions to residual variation, population weights, and labeling costs.

The methods retain their established origins: correlation clustering supplies a relational objective, difference estimation and prediction-powered inference supply residual correction, and Neyman allocation supplies the cost-adjusted allocation principle \citep{bansal2004,sarndal1992,ppi2023}. The contribution concerns their connections, conditions, and selection for semantic measurement. Section~\ref{sec:pilot-allocation} develops an explicit two-stage verification baseline from these principles. A claim of a new estimation or allocation method would require a dedicated comparison with that literature.

Sections~\ref{sec:scope}--\ref{sec:interface} define the evidence scope and framework; Sections~\ref{sec:classification}--\ref{sec:measurement} examine variable construction; Sections~\ref{sec:reliability}--\ref{sec:inference} connect validation to population estimation. Section~\ref{sec:evidence} reviews implementation evidence and uses awaiting evaluation. Code accompanies the constructed examples and corpus evaluation.

\section{Scope and targeted evidence review}\label{sec:scope}

The targeted narrative evidence review supports the framework through retrieval of decision-model research and backward tracing to statistical methods. The literature cutoff is 4 October 2026. TypeSafe documentation and the original Jev papers formed the seed set. Expansion searches included \texttt{Jev System One decision models clustering statistical survey} and \texttt{"Calibrated Decisions at Scale" Jev crash narratives}; additional routes covered \texttt{Clef} and \texttt{OpenAI Decisions API}. Method searches covered prediction-powered inference, design-based supervised learning, annotation reviews, survey allocation, and measurement. The existing search log preserves the recorded query strings and direct-source routes.

Discovery used a general web index and community bibliographies. Source checks used arXiv, ACL Anthology, PMLR, NeurIPS proceedings, publisher records, and official model documentation. Materials were included for an inspectable interface specification, an evaluation with identifiable tasks and comparators, or a statistical method relevant to a framework decision. Announcements contributed interface evidence when their descriptions supported comparison; empirical claims required an inspectable evaluation. Unrelated uses of ``decision'' were excluded. The initial seed set's screening counts and discovery sequence were not recorded, and complete screening totals are unavailable. The log therefore documents actual retrieval routes and inclusion reasons without reconstructing a screening flow.

The retained materials comprise nine empirical or implementation preprints, statistical foundations, two related reviews, and official interface sources. The social-science annotation study uses version~2, revised 23 September 2026; the other eight decision-model preprints use version~1. Existing records document checks of bibliographic metadata and cited numerical claims against the relevant original passages. Mutable documentation retains its access date; the Noul response definition was rechecked on 5 October 2026 during this revision, with the literature cutoff unchanged. On 5 October, the two neighboring reviews were revisited for the contribution comparison, and the dataset and fixed model sources for Section~\ref{sec:real-data} were added as case materials. These checks are recorded separately from the original literature retrieval.

Theoretical and empirical support have different roles. Statistical sources establish methods under stated assumptions; the displayed calculations expose those assumptions in the present setting. Documentation establishes interface properties. Empirical studies support results for identified data, comparators, and model versions. Constructed examples demonstrate consequences of analytical choices and motivate designs for evaluation. Section~\ref{sec:evidence} uses these distinctions to interpret provider evaluations, shared authorship, and benchmark reuse, with detailed disclosures in Appendix~\ref{app:provenance}.

\section{Decision interfaces and statistical objects}\label{sec:interface}

\subsection{An operational definition}\label{sec:definition}

A semantic decision interface has two defining features: an explicit semantic criterion for evaluating the supplied record and a prescribed answer space for the resulting judgment. The criterion may be a natural-language instruction, a rubric, or descriptions attached to candidate answers. This definition includes zero-shot classification and language-model-assisted annotation when their queries specify those elements \citep{yin2019,ziems2024}. It provides a common unit of analysis for task criteria, supplied information, answer spaces, and returned fields across implementations. An interface names an operation; a classifier or language model names a model or learning category. Free-text generation with an unrestricted response space and an embedding-only endpoint fall outside this operational definition. Either can supply inputs to a subsequent judgment with explicit semantic criteria and finite answers. An ordinary trained classifier can implement that judgment when its category definitions and returned fields are specified. The definition therefore supports measurement and verification design across overlapping model categories.

Probability output is an additional property that determines whether probability calibration and proper scoring rules apply. A returned label supports classification-error assessment; a complete probability vector supports further analyses of probabilistic accuracy and decisions under specified losses. Persistent conversational state is an implementation property: the supplied context and its update rule form part of the measurement procedure. A single stateless request can satisfy the same interface definition.

Let $X$ denote an observed record, $s(X)$ the information supplied to a model, $q$ a coding criterion, and $\mathcal C=\{c_1,\ldots,c_K\}$ the answer space. For an interface that exposes a complete probability vector, write
\begin{equation}
 \widehat{\boldsymbol p}(X;q,\mathcal C)
 =\bigl(\widehat p_1,\ldots,\widehat p_K\bigr),\qquad
 \widehat p_k\geq0,\quad \sum_{k=1}^{K}\widehat p_k=1.
 \label{eq:interface}
\end{equation}
The notation suppresses the model version and state construction, both of which must be fixed for a reproducible analysis. The vector records the model's assigned probabilities over the specified answers. Its calibration is assessed relative to an explicitly defined outcome and population.

Jev exposes three primitives. Choice returns an option and a probability distribution over nominal alternatives. Score returns probabilities over ordered rubric levels and a numerical score. Noul returns a number in $[0,1]$, documented as the model-assigned probability of a ``yes'' answer to the supplied question \citep{typesafe-intro,typesafe-score,typesafe-noul}. Its statistical interpretation depends on the question, reference labels, and calibration evidence. An ordered rubric with levels $0,\ldots,K-1$ yields a score based on the probability-weighted level index; the numerical spacing is a property of that coding.

The statistical decision rule is a separate operation. Given actions $a\in\mathcal A$ and a loss $L(a,y)$, a plug-in rule is
\begin{equation}
 \widehat a(X)=\argmin_{a\in\mathcal A}
 \sum_{k=1}^{K} L(a,c_k)\widehat p_k(X).
 \label{eq:action}
\end{equation}
Misclassification costs, a referral cost, and the consequences of merging records can all be represented in the loss. Different research objectives can therefore induce different actions from the same probabilities. Estimating a category proportion, retaining a rare event for review, and assigning a final code are distinct uses of the output.

\subsection{Substantive objects, protocol variables, and verification labels}\label{sec:measurement-layers}

Three measurement levels determine the interpretation of an analysis. The \emph{substantive object} is the event, attribute, or construct that motivates the study. The \emph{protocol-defined variable} $Y^q$ records what the specified coding protocol $q$ would assign, including its evidence rules, category boundaries, and adjudication or rater-sampling procedure. The \emph{observed verification label} $Y^{\mathrm{obs}}$ is obtained when that procedure is implemented with the available evidence and reviewers. The protocol's substantive adequacy and the fidelity of its implementation are separate questions. Application-specific validation and models of observer error provide established approaches to these questions \citep{grimmer2013,dawid1979}.

For an objective event, such as a documented service interruption, $q$ specifies the observation window and evidence required for an event code. Reviewers can disagree because of incomplete records, ambiguous rules, or coding mistakes. Adjudication and comparison with external records assess these discrepancies. If the protocol asks whether a narrative \emph{mentions} an interruption, its prevalence concerns recorded mentions; estimating the prevalence of interruptions additionally requires evidence about how events enter the records.

For a subjective rating, such as perceived severity, $q$ specifies whose judgments are of interest and how they are summarized. A consensus category, an expert-panel decision, and a rating drawn from a defined rater population are different variables. Under the last definition, between-rater variation forms part of the target response distribution; under a consensus definition, individual ratings are inputs to an aggregation rule. Agreement measures must therefore be interpreted against the selected target. The verification record should retain the rater assignment, repeated ratings, and adjudication rule needed to assess that interpretation.

When the substantive object admits a numerical variable $Y_*$, define $\mu_*=\E[Y_*]$, $\mu_q=\E[Y^q]$, and $\mu_{\mathrm{obs}}=\E[Y^{\mathrm{obs}}]$. These means must refer to the same population and statistical units on comparable numerical scales. For a fixed corpus, replace expectations by averages over that corpus. Then
\begin{equation}
 \widehat\mu-\mu_*
 = (\widehat\mu-\mu_{\mathrm{obs}})
  +(\mu_{\mathrm{obs}}-\mu_q)
  +(\mu_q-\mu_*).
 \label{eq:measurement-decomposition}
\end{equation}
This identity separates estimation error for the verification variable, implementation discrepancy, and the difference between the protocol and substantive targets. Sampling design, residual correction, and variance estimation address the first term. Protocol audits, validated reference labels, or an identified measurement model provide evidence about the second. Repeat ratings characterize rating variation and agreement; identifying systematic verification bias additionally depends on reference information or model assumptions. External criteria and validation of the construct and record-generation process address the third. For latent constructs without a defined comparable scale, these relationships remain conceptual until a measurement model supplies that scale. Section~\ref{sec:inference} specifies which target its intervals cover. In the intervening probability and loss formulas, $Y$ denotes the declared reference outcome.

\subsection{Implementation and interface comparison}

Pretraining and subsequent adaptation supply information acquired outside the current study. A natural-language criterion accesses that information and specifies its intended application. A zero-shot evaluation consequently describes the absence of supplied target-task training examples under the evaluated protocol. The origin of a checkpoint's training data remains a separate provenance question. The comparison with a target-trained classifier should state both its training labels and the labels used to assess every method.

TypeSafe describes reinforcement learning for calibrated decisions (RLCD) as its post-training objective \citep{typesafe-primer}. That description identifies a developer objective. Its attainment is an empirical property of predictions on a specified task. \citet{liwagle2026} demonstrate a related implementation route in LLM2Jev: probabilities are read from candidate identifiers in an existing language model, with an optional fine-tuning procedure. Their results motivate separating the effects of the backbone, adaptation, readout, and interface in a comparative study.

The source of a probability also matters. A model can report a number in generated text, expose a normalized candidate likelihood, or return a probability through a dedicated head. These procedures can have different error distributions even when their outputs share a schema. Comparisons should preserve the state, label semantics, and failure-handling policy while recording the readout used by each model. This is particularly relevant when a parser failure otherwise appears as a classification error.

Cloudflare's Clef model card describes a Qwen3.8-27B backbone with a joint schema head and a softmax over each question's options; its input can include text, JSON, images, or video \citep{cloudflare2026}. OpenAI's DevDay announcement describes a Luna-based Decisions API with user-defined questions, finite answers, and text or image context, available in limited preview at announcement \citep{openai2026}. The announcement establishes the intended task interface. Endpoint details and probability fields remain unspecified in that source. Table~\ref{tab:interfaces} therefore compares documented properties without assigning the probability-vector assumption to every implementation.

\begin{table}[!htbp]
\centering\small
\caption{Documented implementations within the operational definition. Probability fields determine the applicable analyses.}\label{tab:interfaces}
\begin{tabularx}{\textwidth}{@{}L{2.25cm}YL{3.25cm}Y@{}}
\toprule
System & Task specification & Documented readout & Evidence status \\
\midrule
Jev & State and Choice, Score, or Noul questions & Per-option or binary probabilities; field-specific confidence & Hosted interface documentation and the retained empirical studies \\
Clef & Multimodal state and compatible typed questions & Joint schema head; per-question softmax & Open weights and developer-run Decision Index evaluation \\
LLM2Jev & Context and candidate identifiers & Candidate-token likelihoods; optional adaptation & Research implementation and backbone-specific experiments \\
OpenAI Decisions API & Text or images; questions with finite answers & Returned selection; probability schema unspecified & Official limited-preview announcement; no evaluation of this API retained \\
\bottomrule
\end{tabularx}
\vspace{0.3em}
\parbox{\textwidth}{\footnotesize Sources: \citet{typesafe-intro,cloudflare2026,liwagle2026,openai2026}. The defining properties are explicit criteria and a prescribed answer space. Zero-shot classification and prompted annotation also qualify when they specify these properties. Probability output and persistent state are implementation attributes; calibration applies to the actual readout and declared reference variable.}
\end{table}

\subsection{From the estimand to an estimation design}\label{sec:framework}

The framework follows six connected decisions. First, specify the target population, statistical unit, and estimand. Second, define $q$, the intended variable $Y^q$, and the procedure for obtaining $Y^{\mathrm{obs}}$. Third, acquire interface outputs with a recorded state construction, model version, and answer space. Fourth, prespecify the error measures and structural conditions to be tested. Fifth, design pilot verification and main-stage sampling allocation within the budget, retaining the relevant inclusion probabilities. Sixth, carry out validation and estimation, reporting uncertainty from sampling and measurement.

\noindent\textit{Illustration: billing-dispute prevalence.}
For a fixed quarterly complaint corpus, protocol $q$ defines an explicit challenge to a billed charge and specifies treatment of refunds, ambiguity, and blinded adjudication. The resulting labels are $Y^{\mathrm{obs}}$. Freeze the model and query, obtain $f_i$ for every record, and define score strata and validation quantities. A probability pilot estimates residual variation and review costs; fresh random samples from the remaining strata supply the main-stage residuals. Equation~\eqref{eq:pilot-estimator} combines both stages to estimate the corpus mean of $Y^{\mathrm{obs}}$. Coding audits and reference evidence address implementation bias; billing records address the separate target of actual billing errors. Section~\ref{sec:real-data} carries out this sampling-design logic on a labeled text corpus.

Table~\ref{tab:roles} relates target quantities to measurement risks and verification designs for four tasks.

\begin{table}[H]
\centering\small
\caption{Estimands, measurement risks, and verification choices.}\label{tab:roles}
\begin{tabularx}{\textwidth}{@{}L{2.15cm}YL{3.7cm}Y@{}}
\toprule
Task & Target quantity and variable construction & Key risk & Verification design \\
\midrule
Classification & Category prevalence from a binary code or class indicator & False-positive and false-negative imbalance; verification-label error & Probability sample under $q$; assess class errors and estimate weighted residuals \\
Clustering & A partition under $q$ and, when meaningful, its group shares; pair judgments supply constraints & Nontransitive local judgments; candidate-edge omissions; ambiguous grouping criterion & Independently coded pairs and a criterion-matched reference partition; audit outside the graph and resample records \\
Ordinal measurement & Mean of a declared score $g(Y^q)$, or prevalence above a defined category boundary & Coding-dependent means; threshold interpretation; rater-population mismatch & Repeated ratings under the chosen protocol; threshold checks and a probability sample for residual correction \\
Record linkage & Entity count or a statistic on linked units; pair judgments supply comparison evidence & Missed candidates; incorrect match prior; false links and global conflicts & Audit candidate and excluded pairs; estimate blocking recall; evaluate the linked structure and downstream quantity \\
\bottomrule
\end{tabularx}
\vspace{0.3em}
\parbox{\textwidth}{\footnotesize Targets refer to the declared protocol. Residual correction estimates means of $Y^{\mathrm{obs}}$ or specified transformations such as $g(Y^{\mathrm{obs}})$. Their relation to $Y^q$ or $g(Y^q)$ requires assessment of implementation discrepancies; the relation of protocol targets to the substantive object requires measurement-validity evidence (Sections~\ref{sec:measurement-layers} and~\ref{sec:inference}).}
\end{table}

\section{Classification and the acquisition of supervision}\label{sec:classification}

\subsection{The role of a coding protocol}

A supervised classifier learns from target labels; a graph-based semi-supervised procedure additionally uses relations among labeled and unlabeled observations \citep{zhu2003}. A semantic decision model can begin with a category description, inclusion rules, and examples supplied in a prompt. Entailment-based zero-shot classification already uses a similar relationship between an observation and a verbalized label \citep{yin2019}. The decision interface standardizes how this relationship is queried and returned.

For mutually exclusive categories, a Choice query provides a natural coding operation. A study with overlapping attributes can evaluate several binary propositions and retain the resulting vector. This vector requires a joint check against the coding protocol: a record can satisfy several substantive attributes, whereas mutually exclusive event states impose constraints on admissible combinations. An ``insufficient information'' category is useful when the protocol distinguishes absence of evidence from evidence of absence. Its use and interpretation should be assessed on adjudicated examples.

The analyst's work consequently includes defining the construct, selecting the state supplied to the model, and designing verification. A model that sees a summary is evaluated on information already filtered by the summarization procedure. If the target is an event described anywhere in the original document, omission during state construction contributes to the measurement error. A useful experiment therefore compares both model readouts and alternative state constructions, while retaining the original record as the unit of sampling.

\subsection{Weak supervision, distillation, and review allocation}

Model judgments can be combined with rules and human labels. Snorkel provides a framework for learning from noisy labeling functions \citep{ratner2017}, while the observer-error model of \citet{dawid1979} supplies an earlier statistical formulation of multiple imperfect annotators. In a decision-model application, rules may identify narrow, verifiable cases, and model queries may cover more varied expressions. The dependence between sources then becomes part of the estimation problem. Rephrased prompts sent to the same checkpoint share training information and may reproduce the same systematic errors.

A downstream classifier can be trained using the returned probability vectors as soft labels. With a sufficiently flexible student and a loss that matches those probabilities, successful fitting reproduces the teacher's conditional outputs. Evidence about correspondence to the study's reference labels comes from a separate labeled sample. A fair cost comparison includes the construction of this sample, teacher queries, student training, and subsequent serving. The relevant advantage may arise from repeated reuse across a large corpus or from rapid adaptation to an infrequent task.

Active learning studies the choice of observations to label \citep{settles2009}. Confidence-based referral is related to this problem when reviewed examples subsequently improve a classifier or its coding rules. For population error assessment, an additional probability sample is useful. Directed review concentrates effort on difficult cases; probability sampling estimates how common those cases are and how frequently errors occur elsewhere. Recording inclusion probabilities permits a common analysis of the two sources of verified labels.

\subsection{Constructed example: accuracy and prevalence}

Consider a finite corpus of 1,000 records containing 200 positive reference labels. Table~\ref{tab:confusion} describes two hypothetical classifiers. Both classify 900 records correctly. Their predicted-positive proportions, $(\mathrm{TP}+\mathrm{FP})/1{,}000$, are 0.20 for A and 0.30 for B; the reference prevalence is 0.20. The example isolates the balance of false positives and false negatives in an uncorrected estimate of prevalence.

\begin{table}[!htbp]
\centering\small
\caption{Constructed classifiers with equal accuracy and different proportion errors.}\label{tab:confusion}
\begin{tabular}{@{}lrrrrrr@{}}
\toprule
Classifier & TP & FN & FP & TN & Accuracy & Predicted-positive proportion \\
\midrule
A & 150 & 50 & 50 & 750 & 0.90 & 0.20 \\
B & 200 & 0 & 100 & 700 & 0.90 & 0.30 \\
\bottomrule
\end{tabular}
\end{table}

For a binary classifier with sensitivity $s$, specificity $c$, and true prevalence $\mu$, its positive rate is
\begin{equation}
 \Prob(\widehat Y=1)=s\mu+(1-c)(1-\mu).
 \label{eq:prevalence}
\end{equation}
This identity underlies classical correction for imperfect screening \citep{rogan1978}. Solving for $\mu$ requires sensitivity and specificity applicable to the target population and a nonzero denominator $s+c-1$. A small denominator amplifies uncertainty in the estimated error rates. Residual correction, developed in Section~\ref{sec:inference}, instead uses a verification sample to estimate the prediction discrepancy directly. Both approaches address the research quantity explicitly, continuing the distinction between document classification and aggregate estimation emphasized by \citet{hopkins2010}.

\section{Clustering through semantic relations}\label{sec:clustering}

\subsection{Geometry, density, and externally supplied relations}

For a representation $h(X_i)$ of $m$ observations, the $k$-means objective is
\begin{equation}
 \min_{z_1,\ldots,z_m,\mu_1,\ldots,\mu_K}
 \sum_{i=1}^{m}\|h(X_i)-\mu_{z_i}\|_2^2,
 \qquad z_i\in\{1,\ldots,K\}.
 \label{eq:kmeans}
\end{equation}
The representation, metric, and number of clusters jointly determine the solution. Lloyd's algorithm supplies the familiar alternating optimization procedure \citep{lloyd1982}. Density-based clustering uses a different structure, illustrated by DBSCAN's neighborhood connectivity \citep{ester1996}. These alternatives express different definitions of a useful group. A corpus of complaints may be grouped by topic, affected product, or event mechanism, each of which can support a different statistical question.

Constraint-based clustering permits prior judgments about which observations belong together \citep{wagstaff2001}. A semantic decision interface can acquire these judgments from a stated criterion. For example, it can compare whether two narratives describe the same failure mechanism, or which of two candidate narratives is more similar to an anchor. In ClusterLLM, triplet judgments guide representation learning and pair judgments guide clustering granularity \citep{clusterllm2023}. This provides a methodological precedent for using language-mediated relations in an otherwise conventional clustering procedure.

These operations support two distinct designs. In a relation-elicitation design, the analyst specifies the grouping criterion and queries local relationships while leaving the partition to an optimization procedure. In a taxonomy-assignment design, the analyst supplies the complete set of categories and asks for membership. The latter has a fixed answer space and can be evaluated as classification. For exploratory clustering, the analysis must specify how the number, interpretation, or resolution of groups is chosen.

\subsection{A soft-constraint formulation}

Let $r_{ij}\in[0,1]$ denote a query's model-assigned probability, or bounded score, for the binary proposition that records $i$ and $j$ should share a group under criterion $q$. When the intended relation is symmetric, query both presentation orders and define $w_{ij}=(r_{ij}+r_{ji})/2$. Equal weighting makes this score invariant to exchanging the record names. On an undirected queried edge set $\mathcal E$, consider
\begin{equation}
 D(\boldsymbol z)=\sum_{(i,j)\in\mathcal E}
 \left[w_{ij}\ind\{z_i\ne z_j\}
 +(1-w_{ij})\ind\{z_i=z_j\}\right].
 \label{eq:correlation}
\end{equation}
Minimizing $D$ is a weighted agreement problem related to correlation clustering \citep{bansal2004}. It separates local evidence from the global requirement that cluster membership form an equivalence relation. Here $w_{ij}$ is a supplied score whose probabilistic interpretation requires validation. Equation~\eqref{eq:correlation} connects the interface to an established objective as an illustrative design. Its evaluation would compare relation sources while holding the candidate graph, optimization rule, and query budget fixed.

\subsection{Uncertainty in symmetrized weights}\label{sec:weight-uncertainty}

The effect of averaging depends on both error variances, their covariance, and their biases. Both query orders must express the same symmetric relation, such as ``same entity''; a directional target such as ``entails'' requires a directed representation.

For a fixed pair, let $Y^q_{ij}$ be the binary relation defined by protocol $q$, and set $\theta_{ij}=\Prob_q(Y^q_{ij}=1\mid X_i,X_j)$. The probability refers to the specified reference mechanism, such as a draw from a defined rater population. With a fixed adjudicated relation, $\theta_{ij}$ is zero or one. Variation in observed verification labels is assessed against this protocol. Separately, across prespecified prompt perturbations or stochastic model evaluations that preserve the relation's meaning, write $r_{ij}=\theta_{ij}+b_{ij}+\epsilon_{ij}$, with $\E[\epsilon_{ij}]=0$, and similarly for the reverse order. Expectations and variances below concern this query randomization, conditional on the pair and protocol. For error variances $\sigma_{ij}^2$, $\sigma_{ji}^2$ and correlation $\rho_{ij}$,
\begin{align}
 \E[w_{ij}]-\theta_{ij}&=\frac{b_{ij}+b_{ji}}{2},\notag\\
 \Var(w_{ij})&=\frac{\sigma_{ij}^2+\sigma_{ji}^2
        +2\rho_{ij}\sigma_{ij}\sigma_{ji}}{4}.
 \label{eq:weight-variance}
\end{align}
With unequal directional variances, the average can have greater variance than the more precise query. With equal variances $\sigma^2$, its variance is $\sigma^2(1+\rho_{ij})/2$, retaining $90\%$ of the single-query variance at correlation $0.8$. Opposite biases can cancel; a shared bias persists. These statements concern variation under the query protocol, with calibration against adjudicated relations assessed separately.

For $R$ repeated scores, assume equal marginal variance $\sigma^2$ and a common pairwise correlation $\rho$ under the same prespecified query-randomization mechanism. When $1+(R-1)\rho>0$,
\begin{equation}
 \Var(\bar r_R)=\frac{\sigma^2}{R}\{1+(R-1)\rho\},\qquad
 R_{\mathrm{eff}}=\frac{R}{1+(R-1)\rho}.
 \label{eq:query-repeats}
\end{equation}
For $R=5$ and $\rho=0.8$, $R_{\mathrm{eff}}=25/21$ and the variance is $0.84\sigma^2$. Correlations are estimated from paired evaluations under a prespecified protocol that preserves the relation's meaning. Repeating an identical deterministic request leaves its score unchanged.

Validate $w_{ij}$ after symmetrization on pairs coded under the declared reference protocol independently of model outputs, including a probability audit outside the candidate graph. Proper scores, calibration curves, and paired order comparisons assess different properties; adjustments use a separate calibration sample. Unequal verification probabilities require sampling weights, and shared records call for record- or source-group resampling. Repeated ratings quantify verification variability separately from repeated model queries.

\subsection{Constructed example: local evidence and a partition}

Let three records have weights $w_{AB}=0.9$, $w_{BC}=0.8$, and $w_{AC}=0.1$. Thresholding at 0.5 accepts the first two links. Transitive closure then joins all three records. Enumerating every partition gives Table~\ref{tab:partitions}; the minimum disagreement occurs at $\{A,B\},\{C\}$. Thus a local threshold followed by connected components solves a different problem from minimizing the stated global loss.

\begin{table}[!htbp]
\centering\small
\caption{Exact costs for a constructed three-record example.}\label{tab:partitions}
\begin{tabular}{@{}lccccc@{}}
\toprule
Partition & $ABC$ & $AB\mid C$ & $AC\mid B$ & $A\mid BC$ & $A\mid B\mid C$ \\
\midrule
$D(\boldsymbol z)$ & 1.2 & 1.0 & 2.6 & 1.2 & 1.8 \\
\bottomrule
\end{tabular}
\end{table}

When $q$ targets co-membership in a common partition, the pair probabilities also satisfy a structural constraint. For any random partition, let $S_{ij}$ denote the event that $i$ and $j$ share a cluster. Transitivity gives $S_{AB}\cap S_{BC}\subseteq S_{AC}$, hence
\begin{equation}
 \Prob(S_{AC})\geq \Prob(S_{AB})+\Prob(S_{BC})-1.
 \label{eq:partition-coherence}
\end{equation}
The constructed scores require a lower bound of 0.7 for $AC$ but assign 0.1. They therefore describe incompatible marginal claims about a partition. Minimizing a partition loss and specifying coherent probabilities over partitions impose distinct requirements. A protocol eliciting pairwise similarity may define a nontransitive relation; the analyst must then specify how that relation supports the intended grouping.

\subsection{Query design and evaluation}

Querying every unordered pair requires $m(m-1)/2$ comparisons, with twice as many calls under the symmetrization above. A practical design can first construct a candidate graph from embeddings, then query a mixture of nearby pairs, current cluster boundaries, and randomly sampled distant pairs. The last component measures relationships missed by the candidate graph. If the purpose is to approximate a full-graph objective, known edge inclusion probabilities can support weighted estimation of that objective; dependence introduced by adaptive querying should be incorporated in the analysis.

Evaluation should share the representation and data split across geometric baselines, human constraints, language-model constraints, and decision-model constraints. Query budgets and any tuning labels should also be reported. Agreement with a reference partition can be measured by an adjusted Rand index \citep{hubert1985}, with the reference constructed for the same criterion $q$. Stability under resampling, alternative descriptions of $q$, and option order assesses additional properties. Resampling documents, and then rebuilding the queried relations, preserves a more appropriate dependence structure than treating all edges as independent observations.

The central opportunity is the elicitation of task-specific relations at a scale that can be compared with manual constraints. Stronger evidence would establish the incremental value of these relations after controlling for representation quality, query selection, and semantic information supplied in the criterion. Such a comparison can determine whether an observed gain arises from better local judgments or from the subsequent optimization procedure.

\section{Measurement, linkage, and related tasks}\label{sec:measurement}

\subsection{Ordinal measurement}

An ordinal protocol orders categories through observable descriptions. A decision model can assign narrative severity under this protocol, after which an ordinal response model relates the coded outcome to covariates \citep{mccullagh1980}. Numerical summaries additionally depend on the scoring convention.

Index the rubric levels by $k=0,\ldots,K-1$. Assigning numerical values $g(k)$ defines a coded outcome $g(Y)$ and a model-based conditional mean
\begin{equation}
 \widehat m_g(X)=\sum_{k=0}^{K-1}g(k)\widehat p_k(X).
 \label{eq:ordinal-mean}
\end{equation}
The coding $g$ defines the score's numerical meaning. The index convention $g(k)=k$ assigns equal increments; interpreting its mean as average severity requires justification for that spacing. Relative to the model-assigned probability distribution, Equation~\eqref{eq:ordinal-mean} is the Bayes action under squared-error loss on the chosen coding. Mean comparisons depend on both the category probabilities and $g$, so order-preserving recodings can alter their magnitude or ranking.

For example, $p^{A}=(1/2,0,1/2)$ and $p^{B}=(0,1,0)$ both have expected index~1. Under $g=(0,1,4)$, their means are 2 and 1, respectively. Their highest-category probabilities are $1/2$ and zero.

The cumulative probabilities and upper-tail probabilities are
\begin{equation}
 \widehat F_j(X)=\sum_{k=0}^{j}\widehat p_k(X),\qquad
 \widehat T_j(X)=1-\widehat F_j(X),\qquad j=0,\ldots,K-2.
 \label{eq:ordinal-cumulative}
\end{equation}
They describe events such as ``at most moderate'' and ``above moderate'', whose meaning follows the category definitions. The identity
\begin{equation}
 \widehat m_g(X)=g(0)+\sum_{j=0}^{K-2}
       \{g(j+1)-g(j)\}\widehat T_j(X)
 \label{eq:ordinal-tail}
\end{equation}
expresses the mean through category increments and tail probabilities. Upper tails at least as large at every threshold give an expected value at least as large under every nondecreasing coding; crossing tails make the comparison coding-dependent. For threshold prevalence, averaging $\widehat T_j(X)$ uses the specified category boundary. Its measurement risks concern boundary interpretation and probability error; verified labels support the residual correction in Section~\ref{sec:inference}. Mean comparisons additionally require a justified scoring convention.

Validation can report threshold-specific reliability and the ranked probability score
\begin{equation}
 \operatorname{RPS}(\widehat{\boldsymbol p},Y)=
 \frac{1}{K-1}\sum_{j=0}^{K-2}
       \bigl\{\widehat F_j-\ind(Y\leq j)\bigr\}^{2},\qquad K\geq2.
 \label{eq:rps}
\end{equation}
This proper score aggregates binary probability errors at equally weighted ordered thresholds \citep{gneiting2007}. Construct validation assesses the rubric's substantive meaning and comparability across groups. For multidimensional rubrics, report component definitions and the aggregation rule.

In a study of Jev~1.13 and three Kev models on 36 ordinal and four nominal datasets, \citet{gao2026} report compressed ordinal scale use. Evaluation should therefore examine category distributions, extreme-category recall, and order sensitivity. Similar aggregate accuracy can coexist with different extreme-category assignments, affecting threshold proportions.

\subsection{Record linkage}

Probabilistic record linkage evaluates evidence that records refer to the same entity \citep{fellegi1969}. Candidate generation comes first: blocking recall, $\Prob(B=1\mid H=1)$, is the fraction of true matches retained as candidates. A probability audit outside the blocks assesses matches excluded before scoring. A semantic model can then supply comparison information for aliases, organizational descriptions, or context-dependent names; TypeSafe provides an entity-alignment example \citep{typesafe-entity}. One-to-one or temporal constraints define the admissible global linkage. These stages contribute distinct errors to entity counts and statistics on linked units.

A returned match probability should be validated within the candidate population created by blocking. This population can have a much higher match prevalence than arbitrary record pairs. Calibration and thresholds consequently depend on how the pairs were generated. For downstream inference, a linkage decision changes both the values attached to a unit and sometimes the number of distinct units. Bayesian entity-resolution approaches represent uncertainty through a joint linkage structure \citep{steorts2016}; they illustrate the additional modeling required to propagate uncertainty beyond isolated pair scores.

Let $H=1$ denote a true match, $B=1$ inclusion by a blocking rule, and $\gamma$ the observed comparison features. Write
\begin{align}
 \pi_B&=\Prob(H=1\mid B=1),\notag\\
 m_B(\gamma)&=\Prob(\gamma\mid H=1,B=1),\qquad
 u_B(\gamma)=\Prob(\gamma\mid H=0,B=1),\notag\\
 \Lambda_B(\gamma)&=\frac{m_B(\gamma)}{u_B(\gamma)},\qquad
 q_B(\gamma)=\frac{\pi_B\Lambda_B(\gamma)}
                     {1-\pi_B+\pi_B\Lambda_B(\gamma)}.
 \label{eq:linkage-posterior}
\end{align}
For continuous features, $m_B$ and $u_B$ are densities with respect to a common measure. These are conditional comparison distributions in the Fellegi--Sunter framework \citep{fellegi1969}. Blocking can change both the prior and the distributions of comparison features, so all three quantities refer to the same candidate population. The joint likelihood ratio accommodates dependence among features; a product of field-specific ratios additionally assumes conditional independence.

Suppose a false link incurs cost $C_{\mathrm{FP}}$ and a missed link costs $C_{\mathrm{FN}}$, with correct actions assigned zero loss. Pairwise expected-loss minimization links when
\begin{equation}
 q_B(\gamma)\geq\frac{C_{\mathrm{FP}}}{C_{\mathrm{FP}}+C_{\mathrm{FN}}}
 \quad\Longleftrightarrow\quad
 \Lambda_B(\gamma)\geq
 \frac{C_{\mathrm{FP}}}{C_{\mathrm{FN}}}\frac{1-\pi_B}{\pi_B}.
 \label{eq:linkage-threshold}
\end{equation}
This rule specifies the pairwise loss before one-to-one or other global linkage constraints are imposed. A clerical-review action would introduce an additional loss and decision region.

For a constructed example, let the match prior among arbitrary pairs be 0.001. Consider a block with prior 0.10, and assume the comparison event has $m=0.80$ and $u=0.02$ in both populations to isolate the prior effect. Its likelihood ratio is~40. Table~\ref{tab:linkage} gives the corresponding posterior and thresholds. With equal error costs, the evidence crosses the linkage threshold inside the block. If a false link costs five times a missed link, the blocked-pair threshold rises to~45 and the same evidence falls below it.

\begin{table}[!htbp]
\centering\small
\caption{Constructed linkage calculation with likelihood ratio 40.}\label{tab:linkage}
\begin{tabularx}{\textwidth}{@{}YL{1.2cm}L{2.2cm}L{2.6cm}L{2.6cm}@{}}
\toprule
Pair population & Prior & Posterior & LR threshold, costs $1{:}1$ & LR threshold, costs $5{:}1$ \\
\midrule
Arbitrary pairs & 0.001 & $40/1039\approx0.0385$ & 999 & 4,995 \\
Blocked candidates & 0.100 & $40/49\approx0.8163$ & 9 & 45 \\
\bottomrule
\end{tabularx}
\end{table}

A model's reported match score can be calibrated directly on sampled blocked pairs, or its comparison features can enter an estimated $m_B/u_B$ model. Applying both a prior correction and a calibration model requires a clear account of which prior is already encoded in the score. Changing the candidate population changes the calibration target and the applicable match prior. Record-level resampling preserves dependence among candidate pairs that share an entity.

\subsection{Ranking and event screening}

Ranking can use relevance scores, pairwise comparisons, or a distribution over a candidate set. \citet{lyu2026} evaluated Jev on three Amazon Reviews domains with candidate sets ranging from 20 to 200 items. They report more gradual latency growth than pointwise Qwen rerankers and higher observed serving latency than recommendation-specific models. These are comparisons of specific serving configurations. Candidate retrieval, total request time, and ranking quality are all relevant to the operational comparison.

Event screening asks whether a record satisfies a definition, such as a reported interruption or an incident requiring investigation. Distance-based outlier detection answers a different question about rarity relative to a representation. An event can be semantically serious and statistically common, or statistically rare and irrelevant to the research criterion. \citet{qin2026} assessed decision readouts on textual evidence from 40 videos and 400 target anchors. Their reported results differ between datasets and readouts. Source videos provide the dependence units for uncertainty calculations; the anchors are repeated measurements within them.

Ranking and screening affect which records reach subsequent analysis. Their statistical evaluation therefore includes the selection rule and the target population represented by the retained records, alongside the local prediction criteria.

\section{Calibration, coherence, and selective prediction}\label{sec:reliability}

\subsection{The meaning of reported confidence}

The numerical field called confidence is defined by its interface. In the TypeSafe documentation accessed on the cutoff date, Choice confidence is
\begin{equation}
 c_{\mathrm{Choice}}=
 \frac{p_{\max}-1/K}{1-1/K},\qquad K>1.
 \label{eq:confidence}
\end{equation}
For a binary choice with $p_{\max}=0.8$, it equals 0.6 \citep{typesafe-confidence}. This transformation describes concentration relative to a uniform distribution. A reliability diagram for the selected option should use the probability whose calibration is being assessed and explicitly identify any transformation. Studies using a verbalized confidence, a top-option probability, and a transformed concentration score assess different quantities.

For a binary probability forecast $p(X)$, calibration is the condition
\begin{equation}
 \E[Y\mid p(X)]=p(X)\quad\text{almost surely}.
 \label{eq:calibration}
\end{equation}
Calibration is a condition on the joint distribution of issued probabilities and outcomes. In multiclass settings, classwise, top-label, and full-vector calibration impose different conditions. A proper scoring rule defines an evaluation loss whose expected value is minimized at the true outcome distribution; strict propriety makes that minimizer unique. The Brier score and logarithmic loss are strictly proper and evaluate probabilistic accuracy, reflecting both calibration and discrimination \citep{gneiting2007}. Reliability diagrams and bin-based calibration errors examine conditional frequencies, with estimates affected by the binning scheme and sample size. Temperature scaling supplies a calibration baseline \citep{guo2017}, fitted separately from final evaluation.

Calibration refers to a specified target population and reference variable. When verification supplies $Y^{\mathrm{obs}}$, Equation~\eqref{eq:calibration} assesses calibration to those observed labels; agreement with $Y^q$ and the substantive object follows the measurement assessment in Section~\ref{sec:measurement-layers}. Changes in the population, candidate set, language, or rubric can alter calibration. A constant prevalence forecast can be calibrated while providing no discrimination. Section~\ref{sec:inference} shows how an informative predictor with a level error can improve estimation after residual correction.

\subsection{Coherence across questions}

Several related queries can express constraints stronger than normalization within each answer. Suppose fine labels partition a broad category $B$. Under a single coherent probability model,
\begin{equation}
 \Prob(Y\in B\mid X)=\sum_{c_k\in B}\Prob(Y=c_k\mid X).
 \label{eq:coherence}
\end{equation}
Asking a separate broad-category question and summing fine-label probabilities provides a diagnostic of this relationship. \citet{joy2026} compares direct, broad-category, and reconstructed predictions for Jev and Laya using 2,500 matched examples per system across three datasets. The study reports discrepancies between the probability paths and changes in prediction quality after hierarchical reconstruction. Thus the sequence of queries is itself an evaluation target.

The partition inequality in \eqref{eq:partition-coherence} is a related constraint for clustering. Logical implication, nested event definitions, and mutually exclusive coding states generate other restrictions. When an analysis combines answers to several questions, it should record these restrictions and assess the joint output. Parallel computation through an interface concerns execution; statistical dependence concerns shared observations, parameters, and errors. Joint covariance estimates should follow the latter structure.

\subsection{Abstention and coverage}

Let $A_\tau(X)$ indicate acceptance by a threshold rule, and suppose its expectation is positive. Classification coverage and selective risk are
\begin{equation}
 \phi(\tau)=\E[A_\tau(X)],\qquad
 R(\tau)=\frac{\E[A_\tau(X)\ind\{\widehat Y\ne Y\}]}{\phi(\tau)}.
 \label{eq:selective}
\end{equation}
Selective classification explicitly studies this tradeoff \citep{geifman2017}. A threshold chosen on a development sample should be evaluated on held-out observations, with uncertainty for both coverage and accepted error. Searching over thresholds on the evaluation outcomes introduces selection into the reported risk. A rigorous risk-control procedure addresses this selection or uses a separate threshold-selection sample.

Referral also changes the composition of the automatically coded records. If high-severity narratives are more ambiguous, acceptance can vary with the outcome. Estimating a population proportion from accepted records then targets a selected population unless the referral mechanism is accounted for. Verification of referred cases, a probability audit of accepted cases, and known selection probabilities can support a complete estimation design.

Prediction sets offer another use of probabilities. Conformal prediction converts conformity scores and an appropriate calibration sample into sets with a coverage guarantee under exchangeability \citep{shafer2008}. Here coverage means inclusion of the true label in a set, whereas $\phi(\tau)$ denotes the fraction of records accepted. Set size and group-specific behavior remain relevant evaluation quantities. Distribution changes affect the assumptions underlying the guarantee and should be included in the study design.

\section{Verification design and population estimation}\label{sec:inference}

\subsection{A common principle and distinct sampling designs}

The estimators in this section target a population mean $\mu=\E[Y]$, including a binary-code prevalence. Here $Y$ abbreviates the observed verification variable $Y^{\mathrm{obs}}$, or a prespecified numerical function of it, such as a category indicator or ordinal score. A model supplies an auxiliary prediction $f(X)$ on that scale. Difference estimation in survey sampling, semi-supervised inference, prediction-powered inference, and design-based supervised learning provide established ways to combine predictions with sampled labels \citep{sarndal1992,zhang2019,ppi2023,egami2023}.

The confidence intervals below target $\mu_{\mathrm{obs}}$, or the corresponding finite-corpus mean, under the stated sampling and fixed-predictor conditions. Their coverage concerns the first term of \eqref{eq:measurement-decomposition}. Coverage for $\mu_q$ or $\mu_*$ requires additional evidence identifying the relevant measurement discrepancies and uncertainty calculations that include their estimation. Repeated ratings estimate variation and sensitivity to adjudication; identifying systematic implementation bias requires a validated reference or identifying assumptions. Random rater assignment belongs in the measurement or sampling model, and shared raters can induce dependence across records. For subjective targets, the specified rater population determines which variation belongs to the estimand.

The common structure is a large collection of auxiliary measurements plus a smaller set of observed discrepancies. Two sampling regimes are especially relevant. In a finite corpus, predictions may be available for every record and labels for a probability sample of those records. In a superpopulation analysis, a labeled sample and an independent unlabeled sample may be drawn from a common distribution. Their variance formulas differ because the first regime includes a known finite-population total and the second estimates a prediction mean from another random sample.

\subsection{An independent-sample difference estimator}

Let $(X_i,Y_i)_{i=1}^{n}$ be independent and identically distributed observations, and let $(\widetilde X_j)_{j=1}^{N}$ be an independent sample from the same marginal distribution of $X$. Condition on a prediction function $f$ fixed independently of these samples. For a fixed coefficient $\lambda$, define
\begin{equation}
 \widehat\mu_\lambda=
 \frac{\lambda}{N}\sum_{j=1}^{N}f(\widetilde X_j)
 +\frac{1}{n}\sum_{i=1}^{n}\{Y_i-\lambda f(X_i)\}.
 \label{eq:ppi}
\end{equation}
Its expectation is $\lambda\E[f(X)]+\E[Y-\lambda f(X)]=\mu$. This cancellation holds for systematic prediction errors as well as for calibrated predictions, provided the two samples represent the specified common population. For finite second moments, independence gives
\begin{equation}
 \Var(\widehat\mu_\lambda)
 =\frac{\lambda^2\Var\{f(X)\}}{N}
  +\frac{\Var\{Y-\lambda f(X)\}}{n}.
 \label{eq:variance-sampling}
\end{equation}
Setting $\lambda=1$ gives the basic mean-estimation form of prediction-powered inference \citep{ppi2023}; $\lambda=0$ gives the labeled sample mean. Write $\sigma_Y^2=\Var(Y)$, $\sigma_f^2=\Var\{f(X)\}$, and $\sigma_{Yf}=\Cov\{Y,f(X)\}$ for the moment comparisons below. The identity locates efficiency in the residual variance, and validity in the sampling and independence conditions.

For fixed $\lambda$, an unbiased variance estimate is the unlabeled sample variance of $\lambda f$ divided by $N$, plus the labeled sample variance of $Y-\lambda f$ divided by $n$. With both sample sizes increasing and a nondegenerate limiting variance, the usual central limit argument yields a Wald interval. Sparse binary outcomes and very small verification samples call for attention to the finite-sample behavior of that approximation. If repeated records share a person or an institution, variance estimation should account for that sampling unit.

Appendix~\ref{app:derivations} gives the moment expansion and variance-minimizing oracle coefficient. The numerical comparison below retains both fixed and optimal weights, so that the contribution of an auxiliary predictor can be assessed against the labeled-only baseline.

\subsection{Constructed example: prediction bias and inferential efficiency}

Let $X$ and $Z$ be independent Bernoulli variables with probability $1/2$, and let
\begin{equation}
 Y\mid X,Z\sim\operatorname{Bernoulli}(0.2+0.6X).
 \label{eq:example-dgp}
\end{equation}
Then $\mu=0.5$ and $\sigma_Y^2=0.25$. Consider three auxiliary predictions: $f_1=0.2+0.6X$, $f_2=0.35+0.6X$, and $f_3=0.2+0.6Z$. Each has variance 0.09. The first two have covariance 0.09 with $Y$; the third has covariance zero. The second prediction overstates the population mean by 0.15, while its centered association with the outcome equals that of the first.

Table~\ref{tab:ppi-example} uses $n=200$ and $N=5{,}000$. The labeled mean has variance 0.00125. The table reports exact population calculations, rounded for presentation. Oracle coefficients describe the information available in the specified distribution; their finite-sample estimation is outside this illustration.

\begin{table}[!htbp]
\centering\small
\caption{Constructed inference example: variance relative to the labeled sample mean.}\label{tab:ppi-example}
\begin{tabularx}{\textwidth}{@{}Yrrrr@{}}
\toprule
Auxiliary prediction & $\E[f]-\mu$ & $\lambda^*$ & $V_1/V_0$ & $V_{\lambda^*}/V_0$ \\
\midrule
$f_1=0.2+0.6X$ & 0.00 & 0.9615 & 0.6544 & 0.6538 \\
$f_2=0.35+0.6X$ & 0.15 & 0.9615 & 0.6544 & 0.6538 \\
$f_3=0.2+0.6Z$ & 0.00 & 0.0000 & 1.3744 & 1.0000 \\
\bottomrule
\end{tabularx}
\vspace{0.3em}
\parbox{\textwidth}{\footnotesize $V_\lambda=\Var(\widehat\mu_\lambda)$ under independent labeled and unlabeled samples. The bias column concerns the uncorrected prediction mean. Every fixed-coefficient corrected estimator in this example is unbiased.}
\end{table}

For $f_2$, the uncorrected prediction average has mean squared error $0.15^2+0.09/5{,}000=0.022518$. Fixed-weight correction with $\lambda=1$ removes the mean discrepancy and gives variance 0.000818. For $f_3$, that same weight increases variance to $1.3744V_0$; the optimal weight $\lambda^*=0$ recovers the labeled mean and its variance $V_0$. The two columns in Table~\ref{tab:ppi-example} preserve this distinction between fixed and optimal weighting. Residual correction addresses the mean error, while inferential efficiency depends on covariance with the outcome and the sample design. An empirical application estimates these moments from verified records.

\subsection{A probability sample from a fixed corpus}

Let $U=\{1,\ldots,M\}$ denote a finite corpus with fixed verification outcomes $y_i$ and predictions $f_i$ for every record. The outcomes are assumed to be defined independently of sample selection; the design calculation conditions on this fixed label vector. If a sample $s$ has known inclusion probabilities $\pi_i>0$, the difference estimator of its finite-population mean is
\begin{equation}
 \widehat\mu_{\mathrm{diff}}=
 \frac{1}{M}\sum_{i\in U}f_i
 +\frac{1}{M}\sum_{i\in s}\frac{y_i-f_i}{\pi_i}.
 \label{eq:design}
\end{equation}
For predictions fixed before sampling, taking expectation over the labeling design replaces $\ind\{i\in s\}/\pi_i$ by one, giving the corpus mean $M^{-1}\sum_i y_i$. Under simple random sampling without replacement of $n$ records,
\begin{equation}
 \Var_s(\widehat\mu_{\mathrm{diff}})
 =\left(1-\frac{n}{M}\right)\frac{S_e^2}{n},\qquad
 S_e^2=\frac{1}{M-1}\sum_{i\in U}(e_i-\bar e_U)^2,
 \quad e_i=y_i-f_i.
 \label{eq:finite-var}
\end{equation}
This is the standard residual-variance principle in model-assisted sampling \citep{sarndal1992}. The full prediction mean is fixed in this design. Appendix~\ref{app:finite-weight} gives the fixed-coefficient extension and its SRS oracle weight, which depend on finite-population residual moments. Designs with unequal probabilities or clustered sampling require their corresponding design variance estimators.

The design can oversample low-confidence records, rare predicted classes, or small subgroups while retaining positive inclusion probabilities across the target corpus. Weighted residuals then correct the auxiliary total; a certainty-reviewed subset has inclusion probability one. Coverage of the full frame is the condition that permits correction of the full corpus mean. When verification labels are generated stochastically during review, uncertainty from that measurement process must accompany the conditional design variance.

\subsection{Verification allocation under a cost budget}\label{sec:allocation}

Verification is allocated here to minimize variance for the specified corpus mean under a cost budget. Population weights determine each stratum's contribution to that mean, residual variation determines how much uncertainty verification can remove, and unit cost determines how many labels the budget can support. The resulting design applies classical cost-adjusted Neyman allocation to prediction residuals \citep{sarndal1992}.

Partition the fixed corpus into $H$ strata with sizes $M_h$, population weights $W_h=M_h/M$, and residual variances $S_{e,h}^{2}$. Strata may combine predicted class, confidence band, and a relevant group. Fix the prediction rule before verification and draw independent simple random samples without replacement, of sizes $n_h$, within strata. With $\bar f_U=M^{-1}\sum_{i\in U}f_i$, the estimator and its design variance are
\begin{align}
 \widehat\mu_{\mathrm{str}}&=\bar f_U+\sum_{h=1}^{H}W_h\bar e_{s_h},\notag\\
 V(\boldsymbol n)&=\sum_{h=1}^{H}W_h^2
       \left(1-\frac{n_h}{M_h}\right)\frac{S_{e,h}^{2}}{n_h}.
 \label{eq:stratified-variance}
\end{align}
The weights remain fixed as the allocation changes. A large constant prediction bias can coexist with small residual variance, allowing precise correction once that stratum is represented.

Let $c_h>0$ be the expected cost of obtaining an adjudicated label in stratum $h$, including any specified repeat review, and let $B_v>0$ be the budget after fixed setup costs. The allocation problem is
\begin{equation}
 \min_{\boldsymbol n}V(\boldsymbol n)
 \quad\text{subject to}\quad
 \sum_h c_h n_h\leq B_v,\qquad n_{h,\min}\leq n_h\leq M_h.
 \label{eq:allocation-problem}
\end{equation}
Because $c_h$ is an expected unit cost, the constraint limits expected total verification expenditure; realized costs can fluctuate with review time.
With positive residual variances and inactive size bounds, the continuous optimum allocates $n_h^*\propto W_hS_{e,h}/\sqrt{c_h}$, with the constant set by the budget. Appendix~\ref{app:derivations} gives the Lagrangian calculation and the exact expression in \eqref{eq:allocation-optimum}.
With active bounds, the solution has the form
$n_h=\min\{M_h,\max(n_{h,\min},aW_hS_{e,h}/\sqrt{c_h})\}$,
where $a$ exhausts the available budget unless every stratum is censused. Integer allocations require a feasible rounding or discrete optimization step. Minimum samples support variance estimation and protect strata whose pilot variance happens to be zero. Budget feasibility is checked before imposing those minima.

For a fixed-cost illustration, let $M=10{,}000$, $(M_1,M_2)=(8{,}000,2{,}000)$, $(S_{e,1},S_{e,2})=(0.10,0.40)$, $(c_1,c_2)=(1,4)$, and $B_v=240$. Equation~\eqref{eq:allocation-optimum} gives $(n_1,n_2)=(80,40)$. Allocation proportional to population size at the same cost gives $(120,30)$. Including the finite-population corrections, their variances are 0.000236 and 0.000262667, respectively, a ratio of $177/197$. These are constructed calculations from specified residual variances. They illustrate why reviewing the least confident cases alone leaves out population weights and verification costs.

For a vector of target quantities, each target generally has its own residual or influence-function variance. A design can minimize a prespecified weighted sum of these variances, or impose precision constraints for priority subgroups. The allocation objective and target weights should be fixed before selecting the design; a single mean-optimal allocation answers one specified estimation problem.

\subsection{Pilot estimates and adaptive allocation}\label{sec:pilot-allocation}

The two-stage design serves as a verification baseline: a pilot estimates residual variances and labeling costs, its verified labels remain in the estimator, and its results determine the main-stage allocation. The design makes the required conditioning and use of pilot labels explicit. \citet{zrnic2024} develop model-assisted adaptive data collection more generally.

Pilot estimates supply the unknown quantities in \eqref{eq:allocation-optimum}. Shrinkage toward pooled variances, minimum sample sizes, and sensitivity analysis across plausible variance profiles address sparse pilots. Confidence may define strata; observed residuals inform their allocations.

Let $P$ be the verified pilot, with $m_h$ records in stratum $h$. Conditional on the pilot and any prediction rule fitted from it, compute $f_i$ for every record and hold that rule fixed. Draw independent main-stage simple random samples $s_h^{(2)}$ of sizes $k_h>0$ from the remaining $M_h-m_h$ records in each nonempty stratum, with $k_h$ chosen from the pilot results. Define
\begin{equation}
 \widehat\mu_{\mathrm{2stage}}=\bar f_U
   +\frac{1}{M}\sum_{i\in P}e_i
   +\sum_h\frac{M_h-m_h}{M}\bar e_{s_h^{(2)}}.
 \label{eq:pilot-estimator}
\end{equation}
Conditional expectation over the main-stage samples replaces each residual mean by the mean of its remaining stratum. Equation~\eqref{eq:pilot-estimator} therefore estimates the full corpus mean without bias, conditional on the pilot and fitted rule. The pilot residuals contribute their observed finite-population total. Main-stage allocation uses the remaining budget, population sizes, and residual-variance estimates; Appendix~\ref{app:derivations} gives its conditional variance.

The conditional argument requires fresh random sampling from the remaining frame and a prediction rule fixed during that stage. For further adaptive waves, the selection rule, probabilities, and prediction versions must be retained, with an estimator and uncertainty calculation matched to that sequential design. In particular, a label-dependent stopping rule changes the distribution of a pooled stratum mean. Recording the sampling stages preserves the information needed for an appropriate analysis.

\subsection{Predictor fitting and model changes}

PPI++ develops data-adaptive prediction weights and more general estimators \citep{ppiplus2024}. The exact calculations here condition on a fixed prediction function and coefficient. A coefficient estimated from a pilot and fixed during fresh main-stage sampling fits the conditional argument of Section~\ref{sec:pilot-allocation}: the pilot labels retain their exact residual contribution, and main-stage variances describe the remaining uncertainty. Section~\ref{sec:real-data} evaluates this construction.

The prediction function includes its version, prompt, rubric, and preprocessing. If these are chosen using the same labels that estimate the residual, the fixed-function calculation needs a revised justification. An independent development sample or an appropriate cross-fitting procedure addresses this reuse. Updating a hosted model during collection creates an additional source of variation: predictions from the updated function should be paired with verification under that same function, or analyzed using explicit version strata.


\subsection{Real-corpus evaluation of a verification design}\label{sec:real-data}

We evaluated five estimation schemes on the UCI SMS Spam Collection \citep{sms-data}. The finite population contains all 5,574 labeled rows of the downloaded archive, including 403 repeated-text occurrences retained as distinct row units. The target is the proportion of archived \texttt{spam} labels: 747/5,574, approximately 0.1340. These supplied labels are the fixed $Y^{\mathrm{obs}}$ for the design calculation. The collection combines reported spam and legitimate messages from several sources; its proportion describes that assembled corpus. The evaluation reuses archived labels and records no new human review times or rating disagreements.

The predictor was \nolinkurl{cross-encoder/nli-MiniLM2-L6-H768}, a natural language inference (NLI) model whose card reports SNLI and MultiNLI training \citep{sms-model}. Its pinned, quantized ONNX artifact evaluated each message against fixed descriptions of legitimate communication and unsolicited promotional or fraudulent spam. A softmax over the two entailment logits supplied $f_i$ for spam. Appendix~\ref{app:reproduction} records the descriptions, revision, and readout. Predictions were computed once without fitting or calibration to the corpus labels and retained for this weight comparison.

At each budget of 200 or 400 labels, 10,000 repetitions compared the label mean under simple random sampling (SRS), the fixed uncorrected prediction mean, unit-weight SRS correction, unit-weight pilot-stratified correction, and a pilot-tuned difference estimator. The two SRS estimators shared each random sample without replacement. The two pilot estimators shared their pilot, allocation, and main-stage samples. Four score-rank strata, with original row order breaking ties, each supplied 20 random pilot labels. This 80-label pilot counted toward the total budget. Unit-weight pilot residual variances were shrunk by 20\% toward their average; the remaining 120 or 320 labels minimized the plug-in conditional variance subject to at least ten main-stage labels per stratum. Unit labeling cost was one.

The added estimator used the design-weighted pilot covariance divided by the known full-corpus score variance, clipped to $[-1,1]$ as specified in \eqref{eq:pilot-weight}. This interval permits either sign of association. The coefficient was held fixed during the main stage and used in every term of \eqref{eq:pilot-estimator}, retaining the pilot's observed residual total. The tuning rule was fixed after inspection of the original results and before the additional simulation; no alternative tuning rules were compared. Allocation continued to use the original unit-weight rule, so the paired pilot comparison isolates coefficient adjustment under that allocation. Full-label diagnostics supplied oracle references and evaluation targets, while the operational coefficient used pilot labels.

Bias and MSE refer to the archived-label proportion. Nominal 95\% normal intervals used sample variances and finite-population corrections; both pilot estimators used their conditional variance in \eqref{eq:pilot-variance}, with the applicable residuals. The fixed uncorrected mean uses zero verification labels and has no design-based interval. Comparisons with SRS assess the complete two-stage design, including stratification, allocation, and pilot cost.

The full-corpus correlation was $\rho_{yf}=-0.0545$. The finite-population SRS oracle in \eqref{eq:finite-weight} was $\lambda_U^*=-0.1463$, yielding a variance ratio of 0.99703 relative to the SRS label mean, a 0.30\% reduction. Restricting that oracle to $[0,1]$ gives zero and the label-mean baseline. These full-label calculations are theoretical SRS references. The operational pilot coefficients had means $-0.1454$ and $-0.1474$ and standard deviations 0.2497 and 0.2503 at budgets 200 and 400, respectively. Their substantial sampling variability accompanies the weak score--label association; the two-stage estimator's efficiency is evaluated under its conditional design.

\begin{table}[!htbp]
\centering\small
\caption{Sampling-design evaluation on the fixed SMS corpus: 10,000 repetitions per verification budget.}\label{tab:sms}
\begin{tabularx}{\textwidth}{@{}rYrrrr@{}}
\toprule
Labels & Estimator & Bias & $\mathrm{MSE}\!\times\!10^4$ & Coverage & Width \\
\midrule
0 & Uncorrected prediction mean & 0.416 & 0.173 & --- & --- \\
200 & SRS labels & 0.032 & 5.568 & 95.10 & 9.256 \\
200 & SRS correction, $\lambda=1$ & 0.016 & 6.604 & 94.44 & 10.061 \\
200 & Pilot correction, $\lambda=1$ & 0.007 & 10.117 & 93.96 & 12.419 \\
200 & Pilot-tuned difference & 0.010 & 9.642 & 93.84 & 12.102 \\
400 & SRS labels & -0.024 & 2.637 & 93.69 & 6.420 \\
400 & SRS correction, $\lambda=1$ & -0.017 & 3.113 & 94.67 & 6.971 \\
400 & Pilot correction, $\lambda=1$ & -0.024 & 3.603 & 94.67 & 7.476 \\
400 & Pilot-tuned difference & -0.023 & 3.413 & 94.73 & 7.287 \\
\bottomrule
\end{tabularx}
\vspace{0.3em}
\parbox{\textwidth}{\footnotesize Labels are total verification budgets. Bias and width are in percentage points; coverage is a percentage. Unit-weight correction fixes $\lambda=1$; pilot tuning estimates $\lambda\in[-1,1]$ from 80 pilot labels, included in both two-stage budgets and estimators. Both pilot designs share all sampled labels and the original allocation. Dashes mark the fixed prediction mean's unavailable design interval. Maximum Monte Carlo standard errors across verified rows are 0.032 percentage points for bias, 0.145 in the displayed MSE units, 0.243 percentage points for coverage, and 0.0144 percentage points for mean width.}
\end{table}

Table~\ref{tab:sms} retains the unit-weight comparisons. SRS correction increased MSE by approximately 18\%, consistent with residual variance 0.13689 versus label variance 0.11608. The four unit-weight stratum residual standard deviations ranged from 0.335 to 0.369, providing modest scope for reallocation. Pilot tuning reduced MSE relative to unit-weight pilot correction by 4.7\% at 200 labels and 5.3\% at 400. The paired MSE reductions were 0.475 and 0.190 in the table's units, with Monte Carlo standard errors 0.046 and 0.016. SRS label means retained the lowest MSE among the four verified estimators; tuned pilot MSE was 73\% and 29\% higher, respectively. Coefficient adjustment improved the original pilot estimator, while coefficient variability and its 80-label pilot remained costs of the evaluated design. The SRS oracle calculation identifies little scope for efficiency improvement from a fixed coefficient for this readout.

The uncorrected prediction mean was 0.13818, with fixed error 0.00416 and squared error below the verified estimators' sampling MSE at these budgets. This aggregate agreement coexisted with weak association at the record level. Observed interval coverage ranged from 93.69\% to 95.10\%, reflecting finite-sample normal approximation error. The method ordering concerns this archived SMS corpus, fixed NLI readout, two budgets, and the specified weight schemes. Coverage targets the archived $Y^{\mathrm{obs}}$ mean; the other discrepancies in \eqref{eq:measurement-decomposition} call for reference validation and evidence about message receipt and corpus construction.

\section{Targeted evidence and implications for the framework}\label{sec:evidence}

\subsection{What current implementations provide}

The documented implementations supply explicit criteria, finite answer spaces, and structured selections. Jev, Clef, and LLM2Jev additionally expose probability readouts that support the checks in Section~\ref{sec:reliability}; the retained OpenAI announcement specifies the intended finite-answer interface and limited-preview status at announcement, with its probability schema unspecified \citep{typesafe-intro,cloudflare2026,liwagle2026,openai2026}. A numerical readout permits probability assessment, while records of its model, state, and criterion permit measurement comparisons across runs.

The theoretical basis comes from the cited measurement, clustering, scoring, and sampling methods. Their mathematical implications are conditional on the reference variable, dependence structure, and sampling design stated near the formulas. The implementation documents establish how inputs and outputs are represented. The empirical studies below examine selected prediction and application properties.

\subsection{What the retained studies test}

Table~\ref{tab:evidence} records research design, model identifiers, comparators, checked results, and limits for the nine retained preprints. Identifiers retain the precision reported by each source. Checks cover the cited original passages and numerical claims; the empirical experiments remain those of the source authors. The annotation study uses version~2 and the other eight use version~1, as recorded in Section~\ref{sec:scope}.

\begingroup\small
\begin{longtable}{@{}L{2.2cm}L{3.5cm}L{5.0cm}L{3.7cm}@{}}
\caption{Nine retained preprints: designs, model comparisons, and scope of the evidence.}\label{tab:evidence}\\
\toprule
Source & Design, models, and comparators & Checked scope and reported finding & Principal limit \\
\midrule
\endfirsthead
\multicolumn{4}{l}{\textit{Table~\thetable\ continued}}\\
\toprule
Source & Design, models, and comparators & Checked scope and reported finding & Principal limit \\
\midrule
\endhead
\bottomrule
\endfoot
\citet{deusser2026} & Broad benchmark; Jev~1.13.0 versus Qwen3.8-27B and Gemma-4-E4B option likelihoods & 37 datasets and 346,009 requests. Under fixed prompts, Jev leads Qwen on 27 datasets and Gemma on all 37; Choice and binary-threshold behavior differ. & Two specified readouts; dataset comparisons address the evaluated task collection. \\
\addlinespace
\citet{ibrahim2026} & Annotation benchmark; Jev~1.13 versus 19 language-model comparators & 18 tasks: 3 pilots and 15 main tasks. Jev trails the per-task best language model on 14 main tasks; median macro-F1 deficit 11.6 points on a 0--100 scale; median inference-cost ratio about 1/44. & Comparator chosen separately by task; request costs exclude the full cost of measurement and verification. \\
\addlinespace
\citet{rafe2026} & Decision-gate benchmark; Jev~1.13.0, open decision models, language models, and trained classifiers & On 600 scored in-scope CLINC-150 items, Jev accuracy is 91.3\%, versus 94.8\% for task-trained embedding logistic regression. & Training-label resources differ; threshold and out-of-scope behavior require their own checks. \\
\addlinespace
\citet{crash2026} & Application study; Jev~1.13.0 compared with coded fields and blinded human judgments & 499,500 Texas crash narratives screened; 195,857 coded; 27 questions; 2,416 human judgments. Reference comparisons examine agreement under the respective coding definitions. & Database fields and reviewed narratives can define different measurement targets. \\
\addlinespace
\citet{gao2026} & Measurement evaluation; Jev~1.13 and Kev~0.8B, 4B, and 9B & 36 ordinal and four nominal tasks. Ordinal scale use is compressed relative to reference labels; candidate order affects assignments. & Rubric distributions and order sensitivity are tested; interpretation depends on the reference scale. \\
\addlinespace
\citet{joy2026} & Coherence evaluation; Jev~1.13.0 and English Laya, direct versus hierarchical queries & 2,500 matched examples per system on TREC, CLINC150, and MASSIVE. Direct and reconstructed probabilities disagree; accuracy and calibration can change in different directions. & Hierarchical label paths are tested; partition coherence is a separate structural target. \\
\addlinespace
\citet{liwagle2026} & Implementation study; Qwen3.5-4B and Qwen3-0.6B, candidate likelihoods with and without fine-tuning & Direct readouts support decisions; adaptation benefits vary by backbone and task. The comparisons distinguish readout from fine-tuning. & Results depend on backbone and task; joint statistical uses require their own validation. \\
\addlinespace
\citet{lyu2026} & Reranking benchmark; \texttt{jev-latest} on 30 September 2026; Qwen2.5-7B-Instruct, SASRec, and DCNv2 & Three Amazon Reviews domains; 20--200 candidates. Jev latency grows more gradually than pointwise Qwen latency and exceeds that of the dedicated recommenders. & Moving model alias; selected candidate sets; hosted and local serving conditions differ. \\
\addlinespace
\citet{qin2026} & Exploratory screening; Jev~1.13.0 versus Qwen3-4B-Instruct-2507 readouts & Textual evidence from 40 videos and 400 target anchors. Rankings and probability quality vary by dataset and readout. & Anchors share videos; uncertainty assessment must reflect the source groups. \\
\end{longtable}
\endgroup

These comparisons answer different questions. The broad benchmark compares two specified option-probability baselines, whereas the annotation benchmark selects the best language model separately for each task. Their comparator rules and task populations explain the different performance statements \citep{deusser2026,ibrahim2026}. The crash study illustrates the measurement distinction in Section~\ref{sec:measurement-layers}: agreement with coded database fields and agreement with reviewed narratives refer to different sources of verification \citep{crash2026}. In the decision-gate study, a threshold selected for approximately 5\% in-scope risk accepts 31.0\% of out-of-scope requests \citep{rafe2026}. Ordinal and hierarchical evaluations similarly identify conditions that local accuracy summaries leave unexamined.

The nine papers contain eight listed author sets: \citet{rafe2026,crash2026} share authors and use different application data. CLINC150 recurs in \citet{deusser2026,rafe2026,joy2026}; the social-science tasks in \citet{rafe2026} follow the collection and replication protocol of \citet{ibrahim2026}; ANLI appears in \citet{deusser2026,gao2026}. Clef's developer-run Decision Index also includes CLINC150, BANKING77, and ANLI. These overlaps support protocol comparisons and create dependence in pooled summaries. The accompanying records identify dataset and protocol reuse; exact item intersections remain uncomputed. Appendix~\ref{app:provenance} records affiliations and disclosures. Clef's internal evaluation retains its provider attribution, and TypeSafe's documented numerical, option-order, and context sensitivities provide further stress-test candidates \citep{cloudflare2026,typesafe-jagged}.

\subsection{Which statistical uses remain designs for evaluation}

Section~\ref{sec:real-data} supplies a fixed-corpus prevalence evaluation: it compares point error, interval behavior, and verification allocation under a single predictor. Evaluation across the reviewed decision models would retain that target and budget while varying the prediction source. Newly collected verification labels would additionally support assessment of coding variation, reference validity, and observed review cost.

For clustering, hold the representation, candidate graph, optimizer, and query budget fixed while varying the source of relations. Held-out coded pairs assess local judgments; a criterion-matched reference partition assesses the resulting grouping. Compare transitive closure with minimization of the stated disagreement loss; empirical performance with Jev remains to be evaluated. For ordinal measurement, compare the chosen numerical summary and threshold proportions under the declared rater and scoring protocol. For linkage, assess candidate recall before interpreting match-score calibration.

Verification allocation can be compared through confidence strata, residual-variance pilots, and observed adjudication costs at a common budget. Sections~\ref{sec:allocation}--\ref{sec:pilot-allocation} give a classical allocation rule and a two-stage baseline for this purpose. Their utility is judged by precision for the selected estimand. A common evaluation of these designs across the retained implementations remains an empirical research task.

Reproducibility requires the model identifier, request date, state construction, full criteria, candidate order, probability readout, confidence definition, and failure policy. Verification records additionally specify sampling probabilities, rater assignment, and adjudication. Repeated queries retain their record identifiers so that uncertainty calculations preserve the statistical units.

The framework covers finite-answer semantic measurement and population means or proportions, with clustering and linkage included as variable-construction problems. The reviewed implementation evidence is early, concentrated on Jev-compatible systems, and partly dependent through authorship and reused benchmarks. The separate SMS case evaluates sampling designs using a fixed open NLI model and archived labels. Constructed calculations demonstrate analytical implications under their stated conditions. The corpus evaluation measures conditional sampling performance; transfer to other corpora and the quality of newly collected verification labels remain application-specific questions. Extensions to causal estimands, longitudinal designs, and general regression parameters fall outside the present treatment.

\section{Conclusion}

Semantic decision interfaces participate in measurement through coding criteria, answer spaces, and model-assigned judgments. The estimand specifies the reference variable; protocol validity and implementation fidelity govern its relationship to the substantive object.

Local prediction quality and downstream validity require separate assessment. Equal accuracy can conceal prevalence bias; transitive closure can conflict with a partition objective; ordinal mean comparisons depend on scoring; and fixed auxiliary weights can increase estimation variance.

The value of auxiliary judgments is assessed through precision for a specified estimand and verification budget. Difference estimation, prediction-powered inference, and classical allocation provide established methods for that assessment. The nine reviewed studies document current capabilities and evaluation conditions. For the archived SMS corpus, fixed NLI readout, and 200/400-label budgets, pilot tuning reduced MSE relative to unit-weight pilot correction. SRS label means retained the lowest MSE among the four verified schemes. The comparison supports evaluating score--label association, coefficient-estimation variability, and pilot cost together when choosing a verification design. Its intervals target the archived verification mean; protocol and substantive validity require their corresponding measurement evidence.

\appendix
\section{Variance and allocation calculations}\label{app:derivations}

\subsection{Independent samples and the oracle weight}

Under the independent-sample and fixed-predictor assumptions of Section~\ref{sec:inference}, expanding the residual variance in \eqref{eq:variance-sampling} gives
\begin{align}
 \Var(\widehat\mu_\lambda)
 &=\frac{\lambda^2\sigma_f^2}{N}
 +\frac{\sigma_Y^2+\lambda^2\sigma_f^2-2\lambda\sigma_{Yf}}{n},
 \label{eq:variance}\\
 \sigma_f^2&=\Var[f(X)],\quad
 \sigma_Y^2=\Var(Y),\quad
 \sigma_{Yf}=\Cov(Y,f(X)).\notag
\end{align}
Differentiating \eqref{eq:variance} gives the oracle coefficient, when $\sigma_f^2>0$,
\begin{equation}
 \lambda^*=\frac{\sigma_{Yf}}{(1+n/N)\sigma_f^2},\qquad
 \Var(\widehat\mu_{\lambda^*})
 =\frac{\sigma_Y^2}{n}
 -\frac{\sigma_{Yf}^2}{n(1+n/N)\sigma_f^2}.
 \label{eq:optimal}
\end{equation}
For a prescribed interval of admissible coefficients, the minimizer is its projection onto that interval. If $f$ is constant, the estimator equals the labeled mean for every coefficient. PPI++ develops data-adaptive power tuning and more general estimation procedures \citep{ppiplus2024}. A fitted coefficient has an estimation error of its own: consistency under appropriate regularity conditions supports the corresponding asymptotic argument. The exact unbiasedness calculation above applies to a fixed coefficient, or conditionally to one selected using independent data.

\subsection{Fixed-corpus SRS and its oracle weight}\label{app:finite-weight}

For the fixed corpus in Section~\ref{sec:inference}, define
$S_{ab}=(M-1)^{-1}\sum_{i\in U}(a_i-\bar a_U)(b_i-\bar b_U)$,
with $S_a^2=S_{aa}$. If $s$ is an SRS without replacement and $\lambda$ is fixed before sampling, then
\begin{align}
 \widehat\mu_{\lambda,U}&=\lambda\bar f_U+\overline{(y-\lambda f)}_s,\notag\\
 \Var_s(\widehat\mu_{\lambda,U})
 &=\left(1-\frac nM\right)\frac{S_{y-\lambda f}^2}{n},\qquad
 \lambda_U^*=\frac{S_{yf}}{S_f^2},\quad S_f^2>0.
 \label{eq:finite-weight}
\end{align}
The residual variance expands as $S_y^2-2\lambda S_{yf}+\lambda^2S_f^2$. The unrestricted oracle variance is the label-mean variance multiplied by $1-\rho_{yf}^2$, where $\rho_{yf}=S_{yf}/(S_yS_f)$ for positive variances. Restricting $\lambda$ to an interval projects $\lambda_U^*$ onto it. The factor $1+n/N$ in \eqref{eq:optimal} belongs to the independent-sample design, whose auxiliary mean is estimated from a random sample. Here $\bar f_U$ is known. These oracle comparisons concern fixed-coefficient SRS; the pilot design uses its own conditional variance in \eqref{eq:pilot-variance}.

\subsection{Cost-adjusted residual allocation}

For fixed strata, the finite-population correction in \eqref{eq:stratified-variance} is constant in the allocation. The interior Lagrangian is $\sum_h W_h^2S_{e,h}^2/n_h+\eta(\sum_h c_hn_h-B_v)$, up to that constant. Its stationarity condition is $W_h^2S_{e,h}^2/n_h^2=\eta c_h$.
With positive residual variances and inactive size bounds, imposing the budget gives
\begin{align}
 n_h^*&=B_v\,
 \frac{W_h S_{e,h}/\sqrt{c_h}}
      {\sum_g W_g S_{e,g}\sqrt{c_g}},\notag\\
 V(\boldsymbol n^*)&=
 \frac{\bigl(\sum_h W_h S_{e,h}\sqrt{c_h}\bigr)^2}{B_v}
       -\sum_h\frac{W_h^2S_{e,h}^{2}}{M_h}.
 \label{eq:allocation-optimum}
\end{align}

The bounds and feasible integer allocation described in Section~\ref{sec:allocation} apply when the interior solution is unavailable.

\subsection{Conditional variance after the pilot}

Conditional expectation over the main-stage samples replaces each residual mean by the mean of its remaining stratum, so \eqref{eq:pilot-estimator} is unbiased for the full corpus mean. The observed pilot residuals contribute their exact finite-population total. Its conditional variance is
\begin{equation}
 \Var(\widehat\mu_{\mathrm{2stage}}\mid P,\boldsymbol y_P)
 =\sum_h\left(\frac{M_h-m_h}{M}\right)^2
     \left(1-\frac{k_h}{M_h-m_h}\right)
     \frac{S_{e,h,\mathrm{rem}}^2}{k_h},
 \label{eq:pilot-variance}
\end{equation}
where $S_{e,h,\mathrm{rem}}^2$ is the finite-population residual variance among the remaining records, conditional on the fitted rule. Empty or fully censused remaining strata contribute zero. The pilot cost is deducted before choosing the main-stage allocation, and the remaining population weights replace $W_h$ in the optimization. Sample residual variances from the main stage estimate the terms in \eqref{eq:pilot-variance} when the sample sizes permit.


For the additional SMS estimator, let $P_h$ contain the $m_h=20$ pilot records in stratum $h$. The score mean and variance are known from the full corpus. The pilot covariance estimate and bounded coefficient are
\begin{equation}
 \widehat S_{yf,P}=\frac{M}{M-1}\sum_h\frac{W_h}{m_h}
       \sum_{i\in P_h}y_i(f_i-\bar f_U),\qquad
 \widehat\lambda_P=\left[\frac{\widehat S_{yf,P}}{S_f^2}\right]_{-1}^{1},
 \label{eq:pilot-weight}
\end{equation}
where $[x]_{-1}^{1}=\max(-1,\min(1,x))$; set the coefficient to zero when $S_f^2=0$. Under the stratified pilot design, $\widehat S_{yf,P}$ is unbiased for $S_{yf}$, since the population sum of $f_i-\bar f_U$ is zero. The bounds permit either sign of association and limit coefficient variability; they were fixed before this extension was simulated, after the original results had been inspected.

Substitute $\widehat\lambda_P f_i$ for $f_i$ throughout \eqref{eq:pilot-estimator}, including its full-corpus mean and pilot residual total. Conditional on the pilot, both the coefficient and main-stage allocation are fixed. Thus the estimator remains conditionally unbiased for the archived-label mean, and \eqref{eq:pilot-variance} applies with residuals $y_i-\widehat\lambda_P f_i$. Main-stage sample variances estimate that conditional variance. Because the conditional mean equals the same corpus target for every pilot, total variance equals the expectation of the conditional variance over pilots. This accounts for coefficient-estimation variability within the actual two-stage design; the normal interval retains its finite-sample approximation error.


\section{Reproduction of examples and the corpus evaluation}\label{app:reproduction}

The accompanying \nolinkurl{scripts/reproduce_examples.py} uses Python's standard library and exact rational arithmetic. It enumerates the confusion matrices in Table~\ref{tab:confusion}, all partitions in Table~\ref{tab:partitions}, and the eight $(X,Z,Y)$ cells in \eqref{eq:example-dgp}. It also checks the correlated-query variances, ordinal score and tail identities, linkage posteriors and thresholds, and the stratified allocation calculation. An eight-record enumeration verifies conditional unbiasedness and the conditional variance of the pilot-based estimator under an outcome-dependent allocation. Exact fractions and decimal values are saved in \texttt{research/worked\_examples.json}. These calculations use constructed inputs and make no calls to model services.

The variance ratios in Table~\ref{tab:ppi-example} are $409/625$ for $f_1$ and $f_2$ at $\lambda=1$, and $859/625$ for $f_3$. Their oracle ratios are $17/26$, $17/26$, and $1$, respectively. These equalities provide simple checks on an independent implementation. All three oracle coefficients belong to $[0,1]$ in this example.

The corpus evaluation is reproduced by \nolinkurl{scripts/sms_design_evaluation.py}. The two hypotheses, in order, were ``This message is legitimate personal or transactional communication.'' and ``This message is unsolicited promotional or fraudulent spam.'' For their entailment logits $\ell_{i0}$ and $\ell_{i1}$, the readout was $f_i=\exp(\ell_{i1})/[\exp(\ell_{i0})+\exp(\ell_{i1})]$, computed with a numerically stable softmax. The fixed repository revision is \nolinkurl{b95119ce93d3e065de6214e38cd4a97b0f2f2c6d}, using \nolinkurl{onnx/model_quint8_avx2.onnx}. Tokenization retained message text, truncated the premise when a pair exceeded 512 tokens, and padded batches of 24 records. No pair exceeded that limit. Execution used the ONNX CPU provider; software versions and artifact hashes are recorded with the predictions.

The stratum sizes were 1,394, 1,394, 1,393, and 1,393. With estimated pilot residual variance $s_h^2$, the allocation used $\widetilde s_h^2=0.8s_h^2+0.2\bar s^2$, where $\bar s^2$ is the average over the four equal-size pilot samples. Starting from ten main-stage labels in each stratum, each remaining label was assigned to the largest reduction in estimated conditional variance; the finite-population correction is constant in that allocation. The weight extension retained this allocation and reused each pilot and main-stage draw for both pilot estimators. NumPy's PCG64 generator used seed 20261005; adding coefficient calculations introduced no further random draws, and the original six verified-result rows were reproduced exactly. Intervals were untruncated at the parameter boundaries. The saved row identifiers, label and score vectors, simulation outputs, and source checksums permit independent reanalysis without rerunning model inference.

\section{Study provenance}\label{app:provenance}

Table~\ref{tab:provenance} summarizes information available in the inspected papers. A listed affiliation describes the publication's author information; a missing disclosure remains unspecified. Detailed source locations, version dates, and benchmark relationships are recorded in the accompanying audit files.

\begingroup\small
\begin{longtable}{@{}L{2.6cm}L{6.0cm}L{6.23cm}@{}}
\caption{Affiliations and study-specific support disclosed in the retained preprints.}\label{tab:provenance}\\
\toprule
Study & Listed affiliation or contact evidence & Support and relationship record \\
\midrule
\endfirsthead
\multicolumn{3}{l}{\textit{Table~\thetable\ continued}}\\
\toprule
Study & Listed affiliation or contact evidence & Support and relationship record \\
\midrule
\endhead
\bottomrule
\endfoot
Deu{\ss}er et al. & University of Bonn; Lamarr Institute; Fraunhofer IAIS & Acknowledged German federal and North-Rhine Westphalia funding through Lamarr. \\
\addlinespace
Ibrahim and Zaki & New York University Abu Dhabi & Study-specific funding and competing-interest disclosure not located. \\
\addlinespace
Rafe and Das, both papers & Texas State University; Ingram School of Engineering in the crash paper & Shared authorship; funding and competing-interest disclosure not located in either paper. \\
\addlinespace
Gao et al. & Jilin University & Study-specific funding and competing-interest disclosure not located. \\
\addlinespace
Joy & Universiti Malaya & Study-specific funding and competing-interest disclosure not located. \\
\addlinespace
Li and Wagle & Microsoft email addresses in the author block & Contact-domain evidence; further study support and competing-interest disclosure unspecified. \\
\addlinespace
Lyu and Xia & Singapore Management University; Meta AI & Industry affiliation recorded; further study support and competing-interest disclosure unspecified. \\
\addlinespace
Qin et al. & Independent researchers, as stated in the author block & Study-specific funding and competing-interest disclosure not located. \\
\end{longtable}
\endgroup

\section{Evidence and preparation records}

The companion search log and evidence matrix record retrieval routes, inclusion reasons, and claim attribution. Preprint, numerical-claim, study-provenance, and dataset-overlap records identify versions, source locations, and the scope of checks. The initial discovery sequence and screening counts are unavailable, as noted in Section~\ref{sec:scope}. These records document AI-assisted comparisons with source texts; empirical results remain attributed to the cited studies.

\section*{Materials and code availability}
The LaTeX source, analysis code, fixed predictions, simulation outputs, and supporting evidence records are archived at \url{https://github.com/LeoAndJellyfish/semantic-decision-interfaces/releases/tag/v1.0.0}. Reproduction instructions and file checksums accompany the archive. UCI distributes the SMS Spam Collection under CC BY 4.0; the model card specifies Apache 2.0. Model weights are retrieved from the pinned revision recorded in Appendix~\ref{app:reproduction}.

\section*{Funding and competing interests}
The author reports no employment, consultancy, investment, or funding relationships with the model providers discussed in this article. This study received no project funding or complimentary API credits.

\section*{Use of generative AI}
OpenAI Codex assisted literature discovery, source cross-checking, English drafting and revision, LaTeX preparation, and development and execution of the reproducibility code. Constructed examples were checked through exact-arithmetic calculations. For the SMS case, Codex implemented and executed the fixed-model predictions, original sampling simulations, and subsequent pilot-weight comparison described in Section~\ref{sec:real-data}; the source corpus supplies the reference labels. Findings from the nine reviewed studies remain attributed to their original authors. The accompanying records identify these distinct sources of results. The author has reviewed the scientific content, references, code, and numerical results and takes responsibility for the manuscript and its declarations.

\begingroup
\small\raggedright
\setlength{\bibsep}{0.4em}
\interlinepenalty=10000
\begin{thebibliography}{99}

\bibitem[Almeida and Hidalgo(2011)]{sms-data}
Almeida, T. and Hidalgo, J. (2011).
\href{https://doi.org/10.24432/C5CC84}{SMS Spam Collection}.
UCI Machine Learning Repository. Dataset, accessed 5 October 2026.


\bibitem[Angelopoulos et~al.(2023)]{ppi2023}
Angelopoulos, A.~N., Bates, S., Fannjiang, C., Jordan, M.~I., and Zrnic, T. (2023).
\href{https://doi.org/10.1126/science.adi6000}{Prediction-powered inference}.
\emph{Science}, 382(6671), 669--674.

\bibitem[Angelopoulos et~al.(2024)]{ppiplus2024}
Angelopoulos, A.~N., Duchi, J.~C., and Zrnic, T. (2024).
\href{https://arxiv.org/abs/2311.01453v2}{PPI++: Efficient prediction-powered inference}.
arXiv:2311.01453v2. Preprint, revised 26 March 2024.

\bibitem[Bansal et~al.(2004)]{bansal2004}
Bansal, N., Blum, A., and Chawla, S. (2004).
\href{https://doi.org/10.1023/B:MACH.0000033116.57574.95}{Correlation clustering}.
\emph{Machine Learning}, 56, 89--113.

\bibitem[Cloudflare(2026)]{cloudflare2026}
Cloudflare (2026).
\href{https://huggingface.co/Cloudflare/clef}{Clef: Model card and Decision Index results}.
Official model documentation and developer-run evaluation, accessed 4 October 2026.

\bibitem[Dawid and Skene(1979)]{dawid1979}
Dawid, A.~P. and Skene, A.~M. (1979).
\href{https://doi.org/10.2307/2346806}{Maximum likelihood estimation of observer error-rates using the EM algorithm}.
\emph{Applied Statistics}, 28(1), 20--28.

\bibitem[Deu{\ss}er et~al.(2026)]{deusser2026}
Deu{\ss}er, T., Sparrenberg, L., and Sifa, R. (2026).
\href{https://arxiv.org/abs/2609.37647v1}{Evaluating and benchmarking the System One model Jev}.
arXiv:2609.37647v1. Preprint.

\bibitem[Egami et~al.(2023)]{egami2023}
Egami, N., Hinck, M., Stewart, B.~M., and Wei, H. (2023).
\href{https://proceedings.neurips.cc/paper_files/paper/2023/hash/d862f7f5445255090de13b825b880d59-Abstract-Conference.html}{Using imperfect surrogates for downstream inference: Design-based supervised learning for social science applications of large language models}.
\emph{Advances in Neural Information Processing Systems}, 36, 68589--68601.

\bibitem[Ester et~al.(1996)]{ester1996}
Ester, M., Kriegel, H.-P., Sander, J., and Xu, X. (1996).
\href{https://f.aaai.org/Papers/KDD/1996/KDD96-037.pdf}{A density-based algorithm for discovering clusters in large spatial databases with noise}.
In \emph{Proceedings of KDD}, pp.~226--231.

\bibitem[Fellegi and Sunter(1969)]{fellegi1969}
Fellegi, I.~P. and Sunter, A.~B. (1969).
\href{https://doi.org/10.1080/01621459.1969.10501049}{A theory for record linkage}.
\emph{Journal of the American Statistical Association}, 64(328), 1183--1210.

\bibitem[Gao et~al.(2026)]{gao2026}
Gao, T., Li, J., Li, Z., Chang, Y., and Wu, Y. (2026).
\href{https://arxiv.org/abs/2609.38827v1}{More choices, fewer decisions: Ordinal-scale bias in JEV-like direct-decision models}.
arXiv:2609.38827v1. Preprint.

\bibitem[Geifman and El-Yaniv(2017)]{geifman2017}
Geifman, Y. and El-Yaniv, R. (2017).
\href{https://proceedings.neurips.cc/paper_files/paper/2017/hash/4a8423d5e91fda00bb7e46540e2b0cf1-Abstract.html}{Selective classification for deep neural networks}.
In \emph{Advances in Neural Information Processing Systems}, 30.

\bibitem[Gneiting and Raftery(2007)]{gneiting2007}
Gneiting, T. and Raftery, A.~E. (2007).
\href{https://doi.org/10.1198/016214506000001437}{Strictly proper scoring rules, prediction, and estimation}.
\emph{Journal of the American Statistical Association}, 102(477), 359--378.

\bibitem[Grimmer and Stewart(2013)]{grimmer2013}
Grimmer, J. and Stewart, B.~M. (2013).
\href{https://doi.org/10.1093/pan/mps028}{Text as data: The promise and pitfalls of automatic content analysis methods for political texts}.
\emph{Political Analysis}, 21(3), 267--297.

\bibitem[Guo et~al.(2017)]{guo2017}
Guo, C., Pleiss, G., Sun, Y., and Weinberger, K.~Q. (2017).
\href{https://proceedings.mlr.press/v70/guo17a.html}{On calibration of modern neural networks}.
\emph{Proceedings of Machine Learning Research}, 70, 1321--1330.

\bibitem[Hopkins and King(2010)]{hopkins2010}
Hopkins, D.~J. and King, G. (2010).
\href{https://doi.org/10.1111/j.1540-5907.2009.00428.x}{A method of automated nonparametric content analysis for social science}.
\emph{American Journal of Political Science}, 54(1), 229--247.

\bibitem[Hubert and Arabie(1985)]{hubert1985}
Hubert, L. and Arabie, P. (1985).
\href{https://doi.org/10.1007/BF01908075}{Comparing partitions}.
\emph{Journal of Classification}, 2, 193--218.

\bibitem[Ibrahim and Zaki(2026)]{ibrahim2026}
Ibrahim, H. and Zaki, Y. (2026).
\href{https://arxiv.org/abs/2609.24574v2}{Evaluating decision models for text annotation in computational social science}.
arXiv:2609.24574v2. Preprint, revised 23 September 2026.

\bibitem[Joy(2026)]{joy2026}
Joy, S.~S. (2026).
\href{https://arxiv.org/abs/2609.33971v1}{Do System One decisions add up? A study of probabilistic coherence}.
arXiv:2609.33971v1. Preprint.

\bibitem[Li and Wagle(2026)]{liwagle2026}
Li, Y. and Wagle, J. (2026).
\href{https://arxiv.org/abs/2610.02076v1}{LLM2Jev: LLMs are already Jev-style decision models---When and how to fine-tune them}.
arXiv:2610.02076v1. Preprint.

\bibitem[Lloyd(1982)]{lloyd1982}
Lloyd, S.~P. (1982).
\href{https://doi.org/10.1109/TIT.1982.1056489}{Least squares quantization in PCM}.
\emph{IEEE Transactions on Information Theory}, 28(2), 129--137.

\bibitem[Lyu and Xia(2026)]{lyu2026}
Lyu, H. and Xia, Y. (2026).
\href{https://arxiv.org/abs/2609.40241v1}{Decision-oriented recommendation reranking: An empirical study of Jev}.
arXiv:2609.40241v1. Preprint.

\bibitem[McCullagh(1980)]{mccullagh1980}
McCullagh, P. (1980).
\href{https://doi.org/10.1111/j.2517-6161.1980.tb01109.x}{Regression models for ordinal data}.
\emph{Journal of the Royal Statistical Society, Series B}, 42(2), 109--127.

\bibitem[OpenAI(2026)]{openai2026}
OpenAI (2026).
\href{https://openai.com/index/devday-2026-recap/}{DevDay 2026 recap}.
Official announcement, 29 September 2026, Decisions API section; accessed 4 October 2026.

\bibitem[Qin et~al.(2026)]{qin2026}
Qin, X., Wang, Y., He, X., Luo, Z., Wu, R., Chen, Y.-S., and Chu, Z. (2026).
\href{https://arxiv.org/abs/2609.34180v1}{Decision readouts for text-mediated video anomaly detection: An exploratory evaluation of Jev and Qwen}.
arXiv:2609.34180v1. Preprint.

\bibitem[Rafe and Das(2026a)]{rafe2026}
Rafe, A. and Das, S. (2026a).
\href{https://arxiv.org/abs/2610.00346v1}{Benchmarking System One decision models against trained classifiers and language models for automated decision gates}.
arXiv:2610.00346v1. Preprint.

\bibitem[Rafe and Das(2026b)]{crash2026}
Rafe, A. and Das, S. (2026b).
\href{https://arxiv.org/abs/2609.24052v1}{Calibrated decisions at scale: Converting police crash narratives into probabilistic crash variables with a System One model (Jev)}.
arXiv:2609.24052v1. Preprint.

\bibitem[Ratner et~al.(2017)]{ratner2017}
Ratner, A., Bach, S.~H., Ehrenberg, H., Fries, J., Wu, S., and R{\'e}, C. (2017).
\href{https://doi.org/10.14778/3157794.3157797}{Snorkel: Rapid training data creation with weak supervision}.
\emph{Proceedings of the VLDB Endowment}, 11(3), 269--282.

\bibitem[Rogan and Gladen(1978)]{rogan1978}
Rogan, W.~J. and Gladen, B. (1978).
\href{https://doi.org/10.1093/oxfordjournals.aje.a112510}{Estimating prevalence from the results of a screening test}.
\emph{American Journal of Epidemiology}, 107(1), 71--76.

\bibitem[S{\"a}rndal et~al.(1992)]{sarndal1992}
S{\"a}rndal, C.-E., Swensson, B., and Wretman, J. (1992).
\href{https://link.springer.com/book/9780387406206}{\emph{Model Assisted Survey Sampling}}.
Springer, New York.

\bibitem[Sentence Transformers(n.d.)]{sms-model}
Sentence Transformers (n.d.).
\href{https://huggingface.co/cross-encoder/nli-MiniLM2-L6-H768}{Cross-encoder for natural language inference: nli-MiniLM2-L6-H768}.
Model card and fixed model files; revision specified in Appendix~\ref{app:reproduction}, accessed 5 October 2026.

\bibitem[Settles(2009)]{settles2009}
Settles, B. (2009).
\href{https://research.cs.wisc.edu/techreports/2009/TR1648.pdf}{Active learning literature survey}.
Computer Sciences Technical Report 1648, University of Wisconsin--Madison.

\bibitem[Shafer and Vovk(2008)]{shafer2008}
Shafer, G. and Vovk, V. (2008).
\href{https://jmlr.org/papers/v9/shafer08a.html}{A tutorial on conformal prediction}.
\emph{Journal of Machine Learning Research}, 9, 371--421.

\bibitem[Steorts et~al.(2016)]{steorts2016}
Steorts, R.~C., Hall, R., and Fienberg, S.~E. (2016).
\href{https://doi.org/10.1080/01621459.2015.1105807}{A Bayesian approach to graphical record linkage and deduplication}.
\emph{Journal of the American Statistical Association}, 111(516), 1660--1672.

\bibitem[Tan et~al.(2024)]{tan2024}
Tan, Z., Li, D., Wang, S., Beigi, A., Jiang, B., Bhattacharjee, A., Karami, M., Li, J., Cheng, L., and Liu, H. (2024).
\href{https://doi.org/10.18653/v1/2024.emnlp-main.54}{Large language models for data annotation and synthesis: A survey}.
In \emph{Proceedings of EMNLP}, pp.~930--957.

\bibitem[TypeSafe AI(2026a)]{typesafe-primer}
TypeSafe AI (2026a).
\href{https://docs.typesafe.ai/introduction/machine-learning-primer}{AI primer}.
Technical documentation, accessed 4 October 2026.

\bibitem[TypeSafe AI(2026b)]{typesafe-confidence}
TypeSafe AI (2026b).
\href{https://docs.typesafe.ai/confidence}{Confidence}.
Technical documentation, accessed 4 October 2026.

\bibitem[TypeSafe AI(2026c)]{typesafe-intro}
TypeSafe AI (2026c).
\href{https://docs.typesafe.ai/introduction}{Introduction}.
Technical documentation, accessed 4 October 2026.

\bibitem[TypeSafe AI(2026d)]{typesafe-jagged}
TypeSafe AI (2026d).
\href{https://docs.typesafe.ai/model-jaggedness/jev-1.13}{Jev 1.13 jaggedness}.
Technical documentation, updated 2 October 2026; accessed 4 October 2026.

\bibitem[TypeSafe AI(2026e)]{typesafe-entity}
TypeSafe AI (2026e).
\href{https://docs.typesafe.ai/cookbooks/entity_alignment}{Knowledge graph entity alignment}.
Technical documentation, accessed 4 October 2026.

\bibitem[TypeSafe AI(2026f)]{typesafe-noul}
TypeSafe AI (2026f).
\href{https://docs.typesafe.ai/primitives/noul}{Noul}.
Technical documentation, accessed 4 October 2026; response definition rechecked 5 October 2026.

\bibitem[TypeSafe AI(2026g)]{typesafe-score}
TypeSafe AI (2026g).
\href{https://docs.typesafe.ai/primitives/score}{Score}.
Technical documentation, accessed 4 October 2026.

\bibitem[Wagstaff et~al.(2001)]{wagstaff2001}
Wagstaff, K., Cardie, C., Rogers, S., and Schroedl, S. (2001).
\href{https://www.cs.cmu.edu/~dgovinda/pdf/icml-2001.pdf}{Constrained k-means clustering with background knowledge}.
In \emph{Proceedings of ICML}, pp.~577--584.

\bibitem[Yin et~al.(2019)]{yin2019}
Yin, W., Hay, J., and Roth, D. (2019).
\href{https://doi.org/10.18653/v1/D19-1404}{Benchmarking zero-shot text classification: Datasets, evaluation and entailment approach}.
In \emph{Proceedings of EMNLP-IJCNLP}, pp.~3914--3923.

\bibitem[Zhang et~al.(2019)]{zhang2019}
Zhang, A., Brown, L.~D., and Cai, T.~T. (2019).
\href{https://doi.org/10.1214/18-AOS1756}{Semi-supervised inference: General theory and estimation of means}.
\emph{The Annals of Statistics}, 47(5), 2538--2566.

\bibitem[Zhang and Gao(2026)]{zhanggao2026}
Zhang, Q. and Gao, C. (2026).
\href{https://doi.org/10.3390/math14122166}{Prediction-powered inference in hybrid measurement regimes: A statistical survey}.
\emph{Mathematics}, 14(12), 2166.

\bibitem[Zhang et~al.(2023)]{clusterllm2023}
Zhang, Y., Wang, Z., and Shang, J. (2023).
\href{https://doi.org/10.18653/v1/2023.emnlp-main.858}{ClusterLLM: Large language models as a guide for text clustering}.
In \emph{Proceedings of EMNLP}, pp.~13903--13920.

\bibitem[Zhu et~al.(2003)]{zhu2003}
Zhu, X., Ghahramani, Z., and Lafferty, J. (2003).
\href{https://mlg.eng.cam.ac.uk/pub/pdf/ZhuGhaLaf03a.pdf}{Semi-supervised learning using Gaussian fields and harmonic functions}.
In \emph{Proceedings of ICML}, pp.~912--919.

\bibitem[Ziems et~al.(2024)]{ziems2024}
Ziems, C., Held, W., Shaikh, O., Chen, J., Zhang, Z., and Yang, D. (2024).
\href{https://doi.org/10.1162/coli_a_00502}{Can large language models transform computational social science?}
\emph{Computational Linguistics}, 50(1), 237--291.

\bibitem[Zrnic and Cand{\`e}s(2024)]{zrnic2024}
Zrnic, T. and Cand{\`e}s, E. (2024).
\href{https://proceedings.mlr.press/v235/zrnic24a.html}{Active statistical inference}.
\emph{Proceedings of Machine Learning Research}, 235, 62993--63010.

\end{thebibliography}
\endgroup
\end{document}
